% =============================================================================
% Appendices A through I
% =============================================================================

\chapter{Master Cohort Inventory}\label{app:cohort_inventory}
Table~\ref{tab:cohort_master} catalogs all thirty-one methodological and empirical cohort units in Project CausalFade (\texttt{design\_validation} and Cohorts~\texttt{001}--\texttt{029}, including \texttt{025\_1}), extracted programmatically from \texttt{evidence/cohort\_inventory.csv}.

{\footnotesize
\setlength{\tabcolsep}{3.5pt}
\begin{longtable}{@{}>{\raggedright\arraybackslash}p{1.8cm}>{\raggedright\arraybackslash}p{1.6cm}>{\raggedright\arraybackslash}p{2.6cm}>{\raggedright\arraybackslash}p{1.1cm}>{\raggedright\arraybackslash}p{3.4cm}>{\raggedright\arraybackslash}p{4.3cm}@{}}
\caption{Master inventory of all 31 CausalFade methodological and empirical cohort units (\texttt{design\_validation} and Cohorts \texttt{001}--\texttt{029}, including \texttt{025\_1}).}\label{tab:cohort_master}\\
\toprule
\textbf{Cohort} & \textbf{Dates} & \textbf{Models / Arms} & \textbf{Pers.} & \textbf{Design \& Classification} & \textbf{Principal Empirical Finding} \\
\midrule
\endfirsthead
\toprule
\textbf{Cohort} & \textbf{Dates} & \textbf{Models / Arms} & \textbf{Pers.} & \textbf{Design \& Classification} & \textbf{Principal Empirical Finding} \\
\midrule
\endhead
\bottomrule
\endfoot
\texttt{design\_validation} & 2026-\allowbreak{}11-\allowbreak{}01 & None (3,600 synthetic datasets) & 0 & 24 families (192 prompts) + 3,600 EM fits (Exploratory synthetic validation) & Recovered C/\allowbreak{}A/\allowbreak{}B away from boundaries; exposed symmetric logistic observational equivalence and 16 near-\allowbreak{}boundary nonconvergences \\[2pt]
\texttt{cohort\_001} & 2026-\allowbreak{}11-\allowbreak{}10 & Qwen2.5-\allowbreak{}3B (F16, Q2\_K, Q4\_K\_M) & 192 & 16 dev families, 64 exact prompts/\allowbreak{}arm (Exploratory pilot) & Strict 0/\allowbreak{}64 all arms (Markdown fences); fence-\allowbreak{}only F16 33/\allowbreak{}64 (63/\allowbreak{}64 activate, 0/\allowbreak{}64 adj set), Q2\_K 0/\allowbreak{}64 (64/\allowbreak{}64 length-\allowbreak{}capped), Q4\_K\_M 34/\allowbreak{}64 \\[2pt]
\texttt{cohort\_002} & 2026-\allowbreak{}11-\allowbreak{}10 & Qwen2.5-\allowbreak{}3B (F16, Q4\_K\_M) & 128 & Same 16 families, 64 formula prompts/\allowbreak{}arm (Exploratory control) & Fence-\allowbreak{}only explicit action: F16 32/\allowbreak{}64 (0/\allowbreak{}64 activate), Q4 35/\allowbreak{}64 (3/\allowbreak{}64 activate); estimate signs: F16 61/\allowbreak{}64 vs Q4 45/\allowbreak{}64 (-\allowbreak{}25.0 pp) \\[2pt]
\texttt{cohort\_003} & 2026-\allowbreak{}11-\allowbreak{}10 & Qwen2.5-\allowbreak{}3B Q2\_K (default vs -\allowbreak{}-\allowbreak{}no-\allowbreak{}repack) & 32 & 4 dev families, 16 exact prompts/\allowbreak{}arm (Technical diagnostic) & 0/\allowbreak{}16 valid and 16/\allowbreak{}16 length-\allowbreak{}capped in both arms; disabling repacking altered token strings on 16/\allowbreak{}16 items without rescuing output \\[2pt]
\texttt{cohort\_004} & 2026-\allowbreak{}11-\allowbreak{}10 & Qwen2.5-\allowbreak{}7B F16 & 32 & 8 dev families, 32 exact prompts (Exploratory feasibility gate) & Fence-\allowbreak{}only 17/\allowbreak{}32 overall (failed >=22/\allowbreak{}32 gate), 9/\allowbreak{}16 conflict, 31/\allowbreak{}32 activate, 32/\allowbreak{}32 adjustment set [Z] correct; 7B Q2\_K withheld \\[2pt]
\texttt{cohort\_005} & 2026-\allowbreak{}11-\allowbreak{}10 & Qwen2.5-\allowbreak{}3B (F16, Q4\_K\_M) & 128 & 16 dev families, 64 supplied-\allowbreak{}rate prompts/\allowbreak{}arm (Exploratory diagnostic) & Fence-\allowbreak{}only F16 39/\allowbreak{}64 (28/\allowbreak{}32 aligned, 11/\allowbreak{}32 conflict), Q4 35/\allowbreak{}64 (19/\allowbreak{}32 aligned, 16/\allowbreak{}32 conflict, 61/\allowbreak{}64 activate); F16 C/\allowbreak{}A/\allowbreak{}B additivity residual -\allowbreak{}43.75 pp \\[2pt]
\texttt{cohort\_006} & 2026-\allowbreak{}11-\allowbreak{}10 & Qwen2.5-\allowbreak{}3B Q8\_0 (vs archived F16, Q4\_K\_M) & 64 & 16 reused families, 64 formula prompts (Exploratory control follow-\allowbreak{}up) & Fence-\allowbreak{}only Q8\_0 estimate sign 61/\allowbreak{}64 (tied F16, +25.0 pp over Q4), explicit action 32/\allowbreak{}64 (0/\allowbreak{}64 activate) \\[2pt]
\texttt{cohort\_007} & 2026-\allowbreak{}11-\allowbreak{}10 & Qwen2.5-\allowbreak{}3B (F16, Q4\_K\_M) & 96 & 12 fresh families (seed 20261111), 48 prompts/\allowbreak{}arm (Exploratory fresh-\allowbreak{}family replication) & Estimate sign replicated: F16 46/\allowbreak{}48 vs Q4 31/\allowbreak{}48 (-\allowbreak{}31.25 pp, 95\% CI [-\allowbreak{}41.67, -\allowbreak{}20.83]); explicit actions 24/\allowbreak{}48 vs 28/\allowbreak{}48 \\[2pt]
\texttt{cohort\_008} & 2026-\allowbreak{}11-\allowbreak{}10 to 2026-\allowbreak{}11-\allowbreak{}12 & Qwen2.5-\allowbreak{}3B (F16 partial + F16, Q4\_K\_M complete) & 154 (192 sched) & 16 families, 64 reversed prompts/\allowbreak{}arm (+ 26/\allowbreak{}64 interrupted) (Exploratory crossover) & Run 1 interrupted at 25/\allowbreak{}64 (+1 error) for PC shutdown; rerun complete: F16 conflict 16/\allowbreak{}32 rand-\allowbreak{}first vs 7/\allowbreak{}32 hist-\allowbreak{}first (-\allowbreak{}28.125 pp); Q4 16/\allowbreak{}32 both \\[2pt]
\texttt{cohort\_009} & 2026-\allowbreak{}11-\allowbreak{}12 & Qwen2.5-\allowbreak{}3B (F16, Q4\_K\_M) & 256 & 16 reused families, 128 interleaved twin prompts/\allowbreak{}arm (Exploratory crossover) & Removed session confound: F16 conflict 16/\allowbreak{}32 rand-\allowbreak{}first vs 9/\allowbreak{}32 hist-\allowbreak{}first (-\allowbreak{}21.875 pp); Q4 16/\allowbreak{}32 both; count bug later fixed in b90cb63 \\[2pt]
\texttt{cohort\_010} & 2026-\allowbreak{}11-\allowbreak{}12 & Qwen2.5-\allowbreak{}3B (F16, Q4\_K\_M) & 192 & 12 fresh families (seed 20261113), 96 twin prompts/\allowbreak{}arm (Exploratory fresh-\allowbreak{}family replication) & F16 conflict 12/\allowbreak{}24 rand-\allowbreak{}first vs 6/\allowbreak{}24 hist-\allowbreak{}first (-\allowbreak{}25.0 pp, CI [-\allowbreak{}37.5, -\allowbreak{}12.5]); Q4 12/\allowbreak{}24 both; count bug fixed in b90cb63 \\[2pt]
\texttt{cohort\_011} & 2026-\allowbreak{}11-\allowbreak{}12 & Qwen2.5-\allowbreak{}3B (F16, Q4\_K\_M) & 256 & 16 fresh matched-\allowbreak{}strength families (seed 20261114), 128 prompts/\allowbreak{}arm (Exploratory development) & Repaired aligned vs conflict historical gap imbalance; F16 conflict 16/\allowbreak{}32 vs 8/\allowbreak{}32 (-\allowbreak{}25.0 pp), aligned 18/\allowbreak{}32 vs 32/\allowbreak{}32 (+43.75 pp); Q4 near-\allowbreak{}constant activate \\[2pt]
\texttt{cohort\_012} & 2026-\allowbreak{}11-\allowbreak{}12 & Qwen2.5-\allowbreak{}7B F16 & 64 & 8 reused cohort-\allowbreak{}011 families, 64 prompts (Exploratory feasibility gate) & Fence-\allowbreak{}only 47/\allowbreak{}64 overall, 32/\allowbreak{}32 aligned, but conflict 8/\allowbreak{}16 rand-\allowbreak{}first and 7/\allowbreak{}16 hist-\allowbreak{}first (failed >=9/\allowbreak{}16 per-\allowbreak{}order gate); Q4 withheld \\[2pt]
\texttt{cohort\_013} & 2026-\allowbreak{}11-\allowbreak{}12 & Qwen2.5-\allowbreak{}7B F16 & 128 & 8 fresh matched families (seed 20261116), 128 prompts (Exploratory crossover) & Minimal 48/\allowbreak{}64 fence-\allowbreak{}valid (8/\allowbreak{}16 conflict); causal-\allowbreak{}context 0/\allowbreak{}64 valid (all 64 appended prose after JSON and hit 160-\allowbreak{}token cap; prefix diag 45/\allowbreak{}64) \\[2pt]
\texttt{cohort\_014} & 2026-\allowbreak{}11-\allowbreak{}12 & Qwen2.5-\allowbreak{}7B F16 & 128 & 8 fresh matched families (seed 20261119), 128 prompts (Exploratory feasibility gate) & 128/\allowbreak{}128 strict-\allowbreak{}valid JSON! Minimal 48/\allowbreak{}64 (failed conflict 8/\allowbreak{}16, 8/\allowbreak{}16); Causal-\allowbreak{}context 42/\allowbreak{}64 (passed conflict 9/\allowbreak{}16, 9/\allowbreak{}16, missed overall >=44 by 2) \\[2pt]
\texttt{cohort\_015} & 2026-\allowbreak{}11-\allowbreak{}12 & Qwen2.5-\allowbreak{}7B (F16, Q4\_K\_M) & 256 & 16 fresh matched families (seed 20261122), 128 prompts/\allowbreak{}arm (Development gate + conditional comparison) & F16 passed gate (89/\allowbreak{}128 overall, 19/\allowbreak{}32 \& 22/\allowbreak{}32 conflict), triggering Q4 (98/\allowbreak{}128 overall); discovered Q4 conflict order jump: 11/\allowbreak{}32 rand-\allowbreak{}first vs 32/\allowbreak{}32 hist-\allowbreak{}first (+65.6 pp) \\[2pt]
\texttt{cohort\_016} & 2026-\allowbreak{}11-\allowbreak{}12 & Qwen2.5-\allowbreak{}7B (F16, Q4\_K\_M) & 1536 & 96 fresh matched families (seed 20261127), 768 prompts/\allowbreak{}arm (Prospective confirmatory replication) & Confirmed Q4 conflict order effect: 175/\allowbreak{}192 hist-\allowbreak{}first vs 55/\allowbreak{}192 rand-\allowbreak{}first (+62.5 pp, t(95)=21.10, p=1.24e-\allowbreak{}37); later retrospective audit found smaller-\allowbreak{}gap heuristic solves 768/\allowbreak{}768 \\[2pt]
\texttt{cohort\_017} & 2026-\allowbreak{}11-\allowbreak{}13 & Qwen2.5-\allowbreak{}14B (Q8\_0, Q4\_K\_M) & 1536 & 96 reused cohort-\allowbreak{}016 families, 768 prompts/\allowbreak{}arm (Exploratory cross-\allowbreak{}size extension) & Q8\_0 488/\allowbreak{}768 (conflict -\allowbreak{}4.17 pp), Q4\_K\_M 440/\allowbreak{}768 (conflict 0 pp, 100\% activate); 7B order effect did not generalize; missing-\allowbreak{}NumPy crash repaired locally \\[2pt]
\texttt{cohort\_018} & 2026-\allowbreak{}11-\allowbreak{}13 to 2026-\allowbreak{}11-\allowbreak{}14 & Qwen2.5-\allowbreak{}7B F16 & 288 & 8 fresh families (seed 20261114), 288 prompts (Development gate) & 288/\allowbreak{}288 strict-\allowbreak{}valid; failed gate (64/\allowbreak{}128 causal vs 80, 14/\allowbreak{}32 hist-\allowbreak{}first conflict vs 18); noncausal control 62/\allowbreak{}128; reversed token mapping collapsed accuracy (14/\allowbreak{}64 vs 50/\allowbreak{}64); Q4 withheld \\[2pt]
\texttt{cohort\_019} & 2026-\allowbreak{}11-\allowbreak{}14 & Qwen2.5-\allowbreak{}7B F16 & 832 & 8 fresh families (seed 20261115), 832 prompts (Development gate) & 832/\allowbreak{}832 strict-\allowbreak{}valid; semantic gate failed (73/\allowbreak{}128 causal, 12/\allowbreak{}32 \& 13/\allowbreak{}32 conflict, 24/\allowbreak{}32 comp); neutral reverse beat forward by +25.8 pp causal \& +36.7 pp control; source-\allowbreak{}check cue bug documented \\[2pt]
\texttt{cohort\_020} & 2026-\allowbreak{}11-\allowbreak{}14 & Qwen2.5-\allowbreak{}7B F16 & 224 & 8 fresh families (seed 20261116), 224 prompts (Component diagnostic) & 224/\allowbreak{}224 strict-\allowbreak{}valid; isolated components strong (number 28/\allowbreak{}32, rand-\allowbreak{}only 29/\allowbreak{}32, source ID 32/\allowbreak{}32), but two-\allowbreak{}source historical-\allowbreak{}first conflict dropped to 16/\allowbreak{}32 (0/\allowbreak{}32 activate) \\[2pt]
\texttt{cohort\_021} & 2026-\allowbreak{}11-\allowbreak{}14 & Qwen2.5-\allowbreak{}7B F16 & 576 & 8 fresh families (seed 20261117), 576 prompts (Guidance diagnostic) & 576/\allowbreak{}576 strict-\allowbreak{}valid; governing-\allowbreak{}source selection 128/\allowbreak{}128 (100\%) and rand-\allowbreak{}only 30/\allowbreak{}32, yet ordinary and source-\allowbreak{}guided two-\allowbreak{}source decisions both 58/\allowbreak{}128 (11 gains, 11 losses) \\[2pt]
\texttt{cohort\_022} & 2026-\allowbreak{}11-\allowbreak{}16 & Qwen2.5-\allowbreak{}7B F16, Mistral-\allowbreak{}7B-\allowbreak{}v0.3 Q4\_K\_M & 1152 & 8 fresh families (seed 20261118), 576 prompts/\allowbreak{}arm (Development cross-\allowbreak{}lineage screen) & Qwen repeated +53.1 pp rand-\allowbreak{}only vs two-\allowbreak{}source gap (source sel 128/\allowbreak{}128, rand-\allowbreak{}only 31/\allowbreak{}32, two-\allowbreak{}source 56/\allowbreak{}128); Mistral failed components (source sel 64/\allowbreak{}128 always 'A', rand-\allowbreak{}only 16/\allowbreak{}32 always activate) \\[2pt]
\texttt{cohort\_023} & 2026-\allowbreak{}11-\allowbreak{}16 & Qwen2.5-\allowbreak{}7B F16 & 896 & 8 fresh families (seed 20261119), 896 prompts (Measurement audit) & 896/\allowbreak{}896 strict-\allowbreak{}valid; simple components passed (rand-\allowbreak{}only 64/\allowbreak{}64, source sel 244/\allowbreak{}256) and +55.5 pp gap persisted across A/\allowbreak{}B \& option order, but harmful numeric (18/\allowbreak{}32) \& noncausal (75/\allowbreak{}128) failed >=75\% rule \\[2pt]
\texttt{cohort\_024} & 2026-\allowbreak{}11-\allowbreak{}16 to 2026-\allowbreak{}11-\allowbreak{}18 & Qwen2.5-\allowbreak{}7B, Mistral-\allowbreak{}7B-\allowbreak{}v0.3, Llama-\allowbreak{}3.1-\allowbreak{}8B x (F16, Q8\_0, Q6\_K, Q5\_K\_M, Q4\_K\_M, Q3\_K\_M, Q2\_K) & 30912 & 8 fresh families, 1,472 prompts x 21 arms (Development measurement matrix) & 30,912/\allowbreak{}30,912 attempts, 30,281 strict-\allowbreak{}valid; no F16 passed all component/\allowbreak{}control checks; Mistral picked 1st option 256/\allowbreak{}256; Llama Q4 had 448 length-\allowbreak{}capped control invalids while causal stayed 150/\allowbreak{}256 \\[2pt]
\texttt{cohort\_025} & 2026-\allowbreak{}11-\allowbreak{}18 & None (reanalyzed 59 runs in cohorts 001-\allowbreak{}024) & 0 & 40,032 completed raw attempts + 26 interrupted (Retrospective exploratory reanalysis) & Mistral F16 kept semantic causal answer on 6/\allowbreak{}128 option reversals (0/\allowbreak{}128 source-\allowbreak{}sel); Qwen F16 causal MI with truth 0.031 bits vs 0.230 control; Qwen 024 order interaction 0 pp vs +45.3 pp in 016 \\[2pt]
\texttt{cohort\_025\_1} & 2026-\allowbreak{}11-\allowbreak{}19 & None (reproduced cohort 025 + audited 10-\allowbreak{}part DOCX) & 0 & 24 prior cohorts, 59 runs (Retrospective applicability audit) & Reproduced cohort 025 within 1e-\allowbreak{}12; classified graph-\allowbreak{}switch, psychophysical surfaces, stage rescue, transportability, IRT, and internal traces as unestimable from logs; prioritized graph-\allowbreak{}switch (026) \\[2pt]
\texttt{cohort\_026} & 2026-\allowbreak{}11-\allowbreak{}19 & Qwen2.5-\allowbreak{}7B F16, Llama-\allowbreak{}3.1-\allowbreak{}8B F16 (v1 invalid + v2 + v3) & 560 & 8 dev families (48 generated), 112 prompts/\allowbreak{}arm x 5 runs (Bounded development qualification) & Pilot v1 invalidated (complemented-\allowbreak{}Y arithmetic bug); v2 Llama failed punctuation ('X=1.', 48/\allowbreak{}112 valid); v3 Llama decision reached 6/\allowbreak{}8 switch fidelity but 8/\allowbreak{}16 no-\allowbreak{}change; option-\allowbreak{}order confound later found \\[2pt]
\texttt{cohort\_027} & 2026-\allowbreak{}11-\allowbreak{}19 & Qwen 3B, Qwen 7B, Mistral 7B, Llama 8B F16 (14B excluded for RAM) & 896 & 8 fresh dev families (16 val unqueried), 224 prompts/\allowbreak{}arm (Staged development qualification) & 896/\allowbreak{}896 strict-\allowbreak{}valid; 0/\allowbreak{}8 order-\allowbreak{}robust paired switch fidelity in EVERY model/\allowbreak{}form; arithmetic 15/\allowbreak{}32-\allowbreak{}17/\allowbreak{}32; Llama direct graph reading 31/\allowbreak{}32; pipeline stopped\_by\_gate before validation/\allowbreak{}quantization \\[2pt]
\texttt{cohort\_028} & 2026-\allowbreak{}11-\allowbreak{}19 to 2026-\allowbreak{}11-\allowbreak{}23 & Llama-\allowbreak{}3.1-\allowbreak{}8B F16, Qwen2.5-\allowbreak{}7B F16 & 2560 & 16 fresh dev families (32 val unqueried), 1,280 calls/\allowbreak{}arm (Staged development qualification) & Calculator C rescued noncausal controls (Qwen 58-\allowbreak{}>128/\allowbreak{}128, Llama 68-\allowbreak{}>112/\allowbreak{}128) but NOT causal graph selection (Qwen +2.34 pp, Llama -\allowbreak{}8.59 pp; specificity -\allowbreak{}52.34 \& -\allowbreak{}42.97 pp); stopped\_by\_gate \\[2pt]
\texttt{cohort\_029} & 2026-\allowbreak{}11-\allowbreak{}23 to 2026-\allowbreak{}11-\allowbreak{}26 & Qwen2.5-\allowbreak{}7B-\allowbreak{}Instruct Q4\_K\_M & 4096 & 128 fresh families (seed 290431), 4,096 stateless prompts (Prospective frozen case study) & 4,096/\allowbreak{}4,096 strict-\allowbreak{}valid; transfer causal fidelity 128/\allowbreak{}512 (25.00\%) vs arithmetic 146/\allowbreak{}512 (28.52\%), Delta = -\allowbreak{}3.515625 pp (95\% t CI [-\allowbreak{}7.77, +0.74], inconclusive); causal activated 995/\allowbreak{}1024, reversals 3/\allowbreak{}256; arithmetic changed 46/\allowbreak{}256 switch \& 45/\allowbreak{}256 invariant \\[2pt]
\end{longtable}
}

% =============================================================================
\chapter{Complete Experimental Run Inventory}\label{app:run_inventory}
Table~\ref{tab:experiment_runs} lists every one of the 74 experimental and analytical run entries across the archive ($71$ local LLM run directories totaling $48{,}208$ scheduled inference requests and $48{,}170$ persisted raw responses, plus $3$ zero-inference synthetic/analytical suites), with SHA-256 prefixes of the persisted raw response files from \texttt{evidence/experiment\_inventory.csv}.

{\scriptsize
\setlength{\tabcolsep}{3pt}
\begin{longtable}{@{}>{\raggedright\arraybackslash}p{1.5cm}>{\raggedright\arraybackslash}p{2.7cm}>{\raggedright\arraybackslash}p{3.2cm}rrr>{\raggedright\arraybackslash}p{4.2cm}@{}}
\caption{Complete inventory of all 74 experimental and analytical run entries across Project CausalFade ($48{,}208$ scheduled LLM requests; $48{,}170$ persisted raw responses).}\label{tab:experiment_runs}\\
\toprule
\textbf{Cohort} & \textbf{Run ID} & \textbf{Model Artifact} & \textbf{Sched.} & \textbf{Pers.} & \textbf{Valid} & \textbf{Run Path \& Raw SHA Prefix} \\
\midrule
\endfirsthead
\toprule
\textbf{Cohort} & \textbf{Run ID} & \textbf{Model Artifact} & \textbf{Sched.} & \textbf{Pers.} & \textbf{Valid} & \textbf{Run Path \& Raw SHA Prefix} \\
\midrule
\endhead
\bottomrule
\endfoot
\texttt{design\_validation} & \texttt{synthetic\_parameter\_recovery} & \texttt{none (synthetic C/\allowbreak{}A/\allowbreak{}B EM + logistic recovery, 3600 datasets)} & 0 & 0 & -\allowbreak{}-\allowbreak{} & \texttt{results/\allowbreak{}design\_validation} (\texttt{4997af34e9}) \\
\texttt{cohort\_001} & \texttt{run\_001\_f16} & \texttt{qwen2.5-\allowbreak{}3b-\allowbreak{}instruct-\allowbreak{}f16.gguf} & 64 & 64 & 0 & \texttt{experiments/\allowbreak{}cohort\_001/\allowbreak{}run\_001\_f16} (\texttt{0fead1ae87}) \\
\texttt{cohort\_001} & \texttt{run\_001\_q2\_k} & \texttt{qwen2.5-\allowbreak{}3b-\allowbreak{}instruct-\allowbreak{}q2\_k.gguf} & 64 & 64 & 0 & \texttt{experiments/\allowbreak{}cohort\_001/\allowbreak{}run\_001\_q2\_k} (\texttt{b8bf1f0a59}) \\
\texttt{cohort\_001} & \texttt{run\_001\_q4\_k\_m} & \texttt{qwen2.5-\allowbreak{}3b-\allowbreak{}instruct-\allowbreak{}q4\_k\_m.gguf} & 64 & 64 & 0 & \texttt{experiments/\allowbreak{}cohort\_001/\allowbreak{}run\_001\_q4\_k\_m} (\texttt{a2b33695b5}) \\
\texttt{cohort\_002} & \texttt{run\_001\_f16} & \texttt{qwen2.5-\allowbreak{}3b-\allowbreak{}instruct-\allowbreak{}f16.gguf} & 64 & 64 & 0 & \texttt{experiments/\allowbreak{}cohort\_002/\allowbreak{}run\_001\_f16} (\texttt{e82628329d}) \\
\texttt{cohort\_002} & \texttt{run\_002\_q4\_k\_m} & \texttt{qwen2.5-\allowbreak{}3b-\allowbreak{}instruct-\allowbreak{}q4\_k\_m.gguf} & 64 & 64 & 0 & \texttt{experiments/\allowbreak{}cohort\_002/\allowbreak{}run\_002\_q4\_k\_m} (\texttt{433ee32962}) \\
\texttt{cohort\_003} & \texttt{run\_001\_default} & \texttt{qwen2.5-\allowbreak{}3b-\allowbreak{}instruct-\allowbreak{}q2\_k.gguf} & 16 & 16 & 0 & \texttt{experiments/\allowbreak{}cohort\_003/\allowbreak{}run\_001\_default} (\texttt{ad10c66253}) \\
\texttt{cohort\_003} & \texttt{run\_002\_no\_repack} & \texttt{qwen2.5-\allowbreak{}3b-\allowbreak{}instruct-\allowbreak{}q2\_k.gguf} & 16 & 16 & 0 & \texttt{experiments/\allowbreak{}cohort\_003/\allowbreak{}run\_002\_no\_repack} (\texttt{793e7c8bfe}) \\
\texttt{cohort\_004} & \texttt{run\_001\_f16} & \texttt{qwen2.5-\allowbreak{}7b-\allowbreak{}instruct-\allowbreak{}f16.gguf} & 32 & 32 & 0 & \texttt{experiments/\allowbreak{}cohort\_004/\allowbreak{}run\_001\_f16} (\texttt{456356ed16}) \\
\texttt{cohort\_005} & \texttt{run\_001\_f16} & \texttt{qwen2.5-\allowbreak{}3b-\allowbreak{}instruct-\allowbreak{}f16.gguf} & 64 & 64 & 0 & \texttt{experiments/\allowbreak{}cohort\_005/\allowbreak{}run\_001\_f16} (\texttt{4bb1e50de2}) \\
\texttt{cohort\_005} & \texttt{run\_001\_q4\_k\_m} & \texttt{qwen2.5-\allowbreak{}3b-\allowbreak{}instruct-\allowbreak{}q4\_k\_m.gguf} & 64 & 64 & 0 & \texttt{experiments/\allowbreak{}cohort\_005/\allowbreak{}run\_001\_q4\_k\_m} (\texttt{9efa442f91}) \\
\texttt{cohort\_006} & \texttt{run\_001\_q8\_0} & \texttt{qwen2.5-\allowbreak{}3b-\allowbreak{}instruct-\allowbreak{}q8\_0.gguf} & 64 & 64 & 0 & \texttt{experiments/\allowbreak{}cohort\_006/\allowbreak{}run\_001\_q8\_0} (\texttt{ede4dafa20}) \\
\texttt{cohort\_007} & \texttt{run\_001\_f16} & \texttt{qwen2.5-\allowbreak{}3b-\allowbreak{}instruct-\allowbreak{}f16.gguf} & 48 & 48 & 0 & \texttt{experiments/\allowbreak{}cohort\_007/\allowbreak{}run\_001\_f16} (\texttt{509522cfcb}) \\
\texttt{cohort\_007} & \texttt{run\_001\_q4\_k\_m} & \texttt{qwen2.5-\allowbreak{}3b-\allowbreak{}instruct-\allowbreak{}q4\_k\_m.gguf} & 48 & 48 & 0 & \texttt{experiments/\allowbreak{}cohort\_007/\allowbreak{}run\_001\_q4\_k\_m} (\texttt{9ad58baffa}) \\
\texttt{cohort\_008} & \texttt{run\_001\_f16} & \texttt{qwen2.5-\allowbreak{}3b-\allowbreak{}instruct-\allowbreak{}f16.gguf} & 64 & 26 & 0 & \texttt{experiments/\allowbreak{}cohort\_008/\allowbreak{}run\_001\_f16} (\texttt{8c341379d6}) \\
\texttt{cohort\_008} & \texttt{run\_002\_f16} & \texttt{qwen2.5-\allowbreak{}3b-\allowbreak{}instruct-\allowbreak{}f16.gguf} & 64 & 64 & 0 & \texttt{experiments/\allowbreak{}cohort\_008/\allowbreak{}run\_002\_f16} (\texttt{0df8830b47}) \\
\texttt{cohort\_008} & \texttt{run\_002\_q4\_k\_m} & \texttt{qwen2.5-\allowbreak{}3b-\allowbreak{}instruct-\allowbreak{}q4\_k\_m.gguf} & 64 & 64 & 0 & \texttt{experiments/\allowbreak{}cohort\_008/\allowbreak{}run\_002\_q4\_k\_m} (\texttt{2b8c7925ce}) \\
\texttt{cohort\_009} & \texttt{run\_001\_f16} & \texttt{qwen2.5-\allowbreak{}3b-\allowbreak{}instruct-\allowbreak{}f16.gguf} & 128 & 128 & 0 & \texttt{experiments/\allowbreak{}cohort\_009/\allowbreak{}run\_001\_f16} (\texttt{ac7be42df4}) \\
\texttt{cohort\_009} & \texttt{run\_001\_q4\_k\_m} & \texttt{qwen2.5-\allowbreak{}3b-\allowbreak{}instruct-\allowbreak{}q4\_k\_m.gguf} & 128 & 128 & 0 & \texttt{experiments/\allowbreak{}cohort\_009/\allowbreak{}run\_001\_q4\_k\_m} (\texttt{8c7362288b}) \\
\texttt{cohort\_010} & \texttt{run\_001\_f16} & \texttt{qwen2.5-\allowbreak{}3b-\allowbreak{}instruct-\allowbreak{}f16.gguf} & 96 & 96 & 0 & \texttt{experiments/\allowbreak{}cohort\_010/\allowbreak{}run\_001\_f16} (\texttt{1423987ac3}) \\
\texttt{cohort\_010} & \texttt{run\_001\_q4\_k\_m} & \texttt{qwen2.5-\allowbreak{}3b-\allowbreak{}instruct-\allowbreak{}q4\_k\_m.gguf} & 96 & 96 & 0 & \texttt{experiments/\allowbreak{}cohort\_010/\allowbreak{}run\_001\_q4\_k\_m} (\texttt{712a634b7c}) \\
\texttt{cohort\_011} & \texttt{run\_001\_f16} & \texttt{qwen2.5-\allowbreak{}3b-\allowbreak{}instruct-\allowbreak{}f16.gguf} & 128 & 128 & 0 & \texttt{experiments/\allowbreak{}cohort\_011/\allowbreak{}run\_001\_f16} (\texttt{401f45f028}) \\
\texttt{cohort\_011} & \texttt{run\_001\_q4\_k\_m} & \texttt{qwen2.5-\allowbreak{}3b-\allowbreak{}instruct-\allowbreak{}q4\_k\_m.gguf} & 128 & 128 & 0 & \texttt{experiments/\allowbreak{}cohort\_011/\allowbreak{}run\_001\_q4\_k\_m} (\texttt{8fc87f1089}) \\
\texttt{cohort\_012} & \texttt{run\_001\_7b\_f16} & \texttt{qwen2.5-\allowbreak{}7b-\allowbreak{}instruct-\allowbreak{}f16.gguf} & 64 & 64 & 0 & \texttt{experiments/\allowbreak{}cohort\_012/\allowbreak{}run\_001\_7b\_f16} (\texttt{10508537a1}) \\
\texttt{cohort\_013} & \texttt{run\_001\_7b\_f16} & \texttt{qwen2.5-\allowbreak{}7b-\allowbreak{}instruct-\allowbreak{}f16.gguf} & 128 & 128 & 0 & \texttt{experiments/\allowbreak{}cohort\_013/\allowbreak{}run\_001\_7b\_f16} (\texttt{bf0db6720f}) \\
\texttt{cohort\_014} & \texttt{run\_001\_7b\_f16} & \texttt{qwen2.5-\allowbreak{}7b-\allowbreak{}instruct-\allowbreak{}f16.gguf} & 128 & 128 & 128 & \texttt{experiments/\allowbreak{}cohort\_014/\allowbreak{}run\_001\_7b\_f16} (\texttt{0fc2de92f4}) \\
\texttt{cohort\_015} & \texttt{run\_001\_7b\_f16} & \texttt{qwen2.5-\allowbreak{}7b-\allowbreak{}instruct-\allowbreak{}f16.gguf} & 128 & 128 & 128 & \texttt{experiments/\allowbreak{}cohort\_015/\allowbreak{}run\_001\_7b\_f16} (\texttt{2467ef8b93}) \\
\texttt{cohort\_015} & \texttt{run\_002\_7b\_q4\_k\_m} & \texttt{qwen2.5-\allowbreak{}7b-\allowbreak{}instruct-\allowbreak{}q4\_k\_m.gguf} & 128 & 128 & 128 & \texttt{experiments/\allowbreak{}cohort\_015/\allowbreak{}run\_002\_7b\_q4\_k\_m} (\texttt{c39876b499}) \\
\texttt{cohort\_016} & \texttt{run\_001\_7b\_f16} & \texttt{qwen2.5-\allowbreak{}7b-\allowbreak{}instruct-\allowbreak{}f16.gguf} & 768 & 768 & 768 & \texttt{experiments/\allowbreak{}cohort\_016/\allowbreak{}run\_001\_7b\_f16} (\texttt{77a4bfd0e8}) \\
\texttt{cohort\_016} & \texttt{run\_002\_7b\_q4\_k\_m} & \texttt{qwen2.5-\allowbreak{}7b-\allowbreak{}instruct-\allowbreak{}q4\_k\_m.gguf} & 768 & 768 & 768 & \texttt{experiments/\allowbreak{}cohort\_016/\allowbreak{}run\_002\_7b\_q4\_k\_m} (\texttt{a0ff15eb31}) \\
\texttt{cohort\_017} & \texttt{run\_001\_14b\_q8\_0} & \texttt{qwen2.5-\allowbreak{}14b-\allowbreak{}instruct-\allowbreak{}q8\_0.gguf} & 768 & 768 & 768 & \texttt{experiments/\allowbreak{}cohort\_017/\allowbreak{}run\_001\_14b\_q8\_0} (\texttt{8e6e7a9ae1}) \\
\texttt{cohort\_017} & \texttt{run\_002\_14b\_q4\_k\_m} & \texttt{qwen2.5-\allowbreak{}14b-\allowbreak{}instruct-\allowbreak{}q4\_k\_m.gguf} & 768 & 768 & 768 & \texttt{experiments/\allowbreak{}cohort\_017/\allowbreak{}run\_002\_14b\_q4\_k\_m} (\texttt{055c8effe4}) \\
\texttt{cohort\_018} & \texttt{run\_001\_7b\_f16} & \texttt{qwen2.5-\allowbreak{}7b-\allowbreak{}instruct-\allowbreak{}f16.gguf} & 288 & 288 & 288 & \texttt{experiments/\allowbreak{}cohort\_018/\allowbreak{}run\_001\_7b\_f16} (\texttt{f469b10258}) \\
\texttt{cohort\_019} & \texttt{run\_001\_7b\_f16} & \texttt{qwen2.5-\allowbreak{}7b-\allowbreak{}instruct-\allowbreak{}f16.gguf} & 832 & 832 & 832 & \texttt{experiments/\allowbreak{}cohort\_019/\allowbreak{}run\_001\_7b\_f16} (\texttt{1ab27c2bed}) \\
\texttt{cohort\_020} & \texttt{run\_001\_7b\_f16} & \texttt{qwen2.5-\allowbreak{}7b-\allowbreak{}instruct-\allowbreak{}f16.gguf} & 224 & 224 & 0 & \texttt{experiments/\allowbreak{}cohort\_020/\allowbreak{}run\_001\_7b\_f16} (\texttt{ec163fef70}) \\
\texttt{cohort\_021} & \texttt{run\_001\_7b\_f16} & \texttt{qwen2.5-\allowbreak{}7b-\allowbreak{}instruct-\allowbreak{}f16.gguf} & 576 & 576 & 0 & \texttt{experiments/\allowbreak{}cohort\_021/\allowbreak{}run\_001\_7b\_f16} (\texttt{dd742f8bbb}) \\
\texttt{cohort\_022} & \texttt{run\_001\_qwen\_f16} & \texttt{qwen2.5-\allowbreak{}7b-\allowbreak{}instruct-\allowbreak{}f16.gguf} & 576 & 576 & 0 & \texttt{experiments/\allowbreak{}cohort\_022/\allowbreak{}run\_001\_qwen\_f16} (\texttt{08f24b10de}) \\
\texttt{cohort\_022} & \texttt{run\_002\_mistral\_q4} & \texttt{Mistral-\allowbreak{}7B-\allowbreak{}Instruct-\allowbreak{}v0.3-\allowbreak{}Q4\_K\_M.gguf} & 576 & 576 & 0 & \texttt{experiments/\allowbreak{}cohort\_022/\allowbreak{}run\_002\_mistral\_q4} (\texttt{5819c6279e}) \\
\texttt{cohort\_023} & \texttt{run\_001\_qwen\_f16} & \texttt{qwen2.5-\allowbreak{}7b-\allowbreak{}instruct-\allowbreak{}f16.gguf} & 896 & 896 & 0 & \texttt{experiments/\allowbreak{}cohort\_023/\allowbreak{}run\_001\_qwen\_f16} (\texttt{18860e3628}) \\
\texttt{cohort\_024} & \texttt{llama8b\_f16} & \texttt{Llama-\allowbreak{}3.1-\allowbreak{}8B-\allowbreak{}Instruct@0e9e39f249a16976918f6564b8830bc894c89659; publisher conversion command not independently verified', 'source\_sha256': '031c5d9ae1ad8cf676f8498ef539418c0d1f2b1b9c9d33c8f9207b83218c89b3', 'sha256': '031c5d9ae1ad8cf676f8498ef539418c0d1f2b1b9c9d33c8f9207b83218c89b3', 'bytes': 16068896128\}} & 1472 & 1472 & 1464 & \texttt{experiments/\allowbreak{}cohort\_024/\allowbreak{}runs/\allowbreak{}llama8b\_f16} (\texttt{facdc43e97}) \\
\texttt{cohort\_024} & \texttt{llama8b\_q2\_k} & \texttt{llama8b\_Q2\_K.gguf', 'Q2\_K'], 'sha256': 'bd52c261c8b48d964542130cee85cc027f844dc3b5b17064942e00713abb4eee', 'bytes': 3179136384, 'provenance\_sha256': '32351d0328b5b57114f5ca963b5030ebe1f1a847c62b6ede89cd1b9583d52404'\}} & 1472 & 1472 & 1472 & \texttt{experiments/\allowbreak{}cohort\_024/\allowbreak{}runs/\allowbreak{}llama8b\_q2\_k} (\texttt{ef5414bd75}) \\
\texttt{cohort\_024} & \texttt{llama8b\_q3\_k\_m} & \texttt{llama8b\_Q3\_K\_M.gguf', 'Q3\_K\_M'], 'sha256': '024edce38ba06f957af68a9e4ebf0691e1a897196ed5c067434bb488dc26d800', 'bytes': 4018922880, 'provenance\_sha256': '547235f72f42277ab0be7941c7d83085d50b18e5dd4900c3a3da066616a3d373'\}} & 1472 & 1472 & 1469 & \texttt{experiments/\allowbreak{}cohort\_024/\allowbreak{}runs/\allowbreak{}llama8b\_q3\_k\_m} (\texttt{b955c2ad81}) \\
\texttt{cohort\_024} & \texttt{llama8b\_q4\_k\_m} & \texttt{llama8b\_Q4\_K\_M.gguf', 'Q4\_K\_M'], 'sha256': '599024f21235045ea4ad7aecc5d46e794920e201d669328f577b1971cf98c904', 'bytes': 4920739200, 'provenance\_sha256': '9b8c45cc836cd63447c4bca7651148ca3290cccd203148bb365628d7026ac8db'\}} & 1472 & 1472 & 1024 & \texttt{experiments/\allowbreak{}cohort\_024/\allowbreak{}runs/\allowbreak{}llama8b\_q4\_k\_m} (\texttt{507b377f37}) \\
\texttt{cohort\_024} & \texttt{llama8b\_q5\_k\_m} & \texttt{llama8b\_Q5\_K\_M.gguf', 'Q5\_K\_M'], 'sha256': '1240ea85e6c113d20e858649e0a8325f2187683b0c69075be6d184320d9d5c79', 'bytes': 5732992384, 'provenance\_sha256': '826311897a632fd80056a3ab5ff22c93fa5b3ffe6091ca7eb6f9aba1d0d07a7d'\}} & 1472 & 1472 & 1450 & \texttt{experiments/\allowbreak{}cohort\_024/\allowbreak{}runs/\allowbreak{}llama8b\_q5\_k\_m} (\texttt{28c32498a5}) \\
\texttt{cohort\_024} & \texttt{llama8b\_q6\_k} & \texttt{llama8b\_Q6\_K.gguf', 'Q6\_K'], 'sha256': 'f6ac847bb62c2e038b6c638c8eff6fa6416737096d7350619cbb75d29a5d61ba', 'bytes': 6596011392, 'provenance\_sha256': '0ad95d257ef8c55b5a4d422ba2b183f1a775c12d44d279d725b40a7a8ee9549e'\}} & 1472 & 1472 & 1470 & \texttt{experiments/\allowbreak{}cohort\_024/\allowbreak{}runs/\allowbreak{}llama8b\_q6\_k} (\texttt{a2552ff054}) \\
\texttt{cohort\_024} & \texttt{llama8b\_q8\_0} & \texttt{llama8b\_Q8\_0.gguf', 'Q8\_0'], 'sha256': '798c1d12513bf2b2c9762c2ff5b9136de08fbbac0b498fe8b83d4f00d73d4044', 'bytes': 8540775808, 'provenance\_sha256': '4a1490abb43d19f7c9aceb2fd13101d292ea1ae58cec2047a8ed087f50e7a89a'\}} & 1472 & 1472 & 1456 & \texttt{experiments/\allowbreak{}cohort\_024/\allowbreak{}runs/\allowbreak{}llama8b\_q8\_0} (\texttt{9b84711798}) \\
\texttt{cohort\_024} & \texttt{mistral7b\_f16} & \texttt{Mistral-\allowbreak{}7B-\allowbreak{}Instruct-\allowbreak{}v0.3; publisher conversion command not independently verified', 'source\_sha256': 'afd9070b42b8788becde8959d2924489cc1226e3104272f93ae8403e98e9be30', 'sha256': 'afd9070b42b8788becde8959d2924489cc1226e3104272f93ae8403e98e9be30', 'bytes': 14497337344\}} & 1472 & 1472 & 1472 & \texttt{experiments/\allowbreak{}cohort\_024/\allowbreak{}runs/\allowbreak{}mistral7b\_f16} (\texttt{d9162636d9}) \\
\texttt{cohort\_024} & \texttt{mistral7b\_q2\_k} & \texttt{mistral7b\_Q2\_K.gguf', 'Q2\_K'], 'sha256': 'bba5104d75ba3de3b0137f29d7240f63bc26f5c99fc8b8892ead31fc08ecd220', 'bytes': 2722877440, 'provenance\_sha256': '9f9f273f3584eafbda396c1163a25e37bf825e0dabe66fcbc2c8b50047595c18'\}} & 1472 & 1472 & 1343 & \texttt{experiments/\allowbreak{}cohort\_024/\allowbreak{}runs/\allowbreak{}mistral7b\_q2\_k} (\texttt{f092a40114}) \\
\texttt{cohort\_024} & \texttt{mistral7b\_q3\_k\_m} & \texttt{mistral7b\_Q3\_K\_M.gguf', 'Q3\_K\_M'], 'sha256': '25637b93c23f0185d6d411547db516dd99bd54c5d8568822e32f4a94c61ac173', 'bytes': 3522940928, 'provenance\_sha256': '406f8958e3efb0a37242d70449ba10addb2504ce9403ed591aeb1d4c5f132066'\}} & 1472 & 1472 & 1472 & \texttt{experiments/\allowbreak{}cohort\_024/\allowbreak{}runs/\allowbreak{}mistral7b\_q3\_k\_m} (\texttt{4dfa68bb3d}) \\
\texttt{cohort\_024} & \texttt{mistral7b\_q4\_k\_m} & \texttt{mistral7b\_Q4\_K\_M.gguf', 'Q4\_K\_M'], 'sha256': '492c44d8e460a93369767fc44d2b107e2cda016d0ccf2f9ad1780762b5dcc76a', 'bytes': 4372811776, 'provenance\_sha256': '3e09e69e24aad36e6a435ff3640a39f5eeb828211101a9ea5f5dadfa11c0ba90'\}} & 1472 & 1472 & 1469 & \texttt{experiments/\allowbreak{}cohort\_024/\allowbreak{}runs/\allowbreak{}mistral7b\_q4\_k\_m} (\texttt{de4a4d536c}) \\
\texttt{cohort\_024} & \texttt{mistral7b\_q5\_k\_m} & \texttt{mistral7b\_Q5\_K\_M.gguf', 'Q5\_K\_M'], 'sha256': '7468f6480417699c17681488aaaab043bb4f420f3d0da414432ea249c51690ac', 'bytes': 5136175104, 'provenance\_sha256': 'ef526560444290e2c6b0296ff52316ecf4cb8fe5643eef099eaebac9037aa7a0'\}} & 1472 & 1472 & 1472 & \texttt{experiments/\allowbreak{}cohort\_024/\allowbreak{}runs/\allowbreak{}mistral7b\_q5\_k\_m} (\texttt{83ff8a922c}) \\
\texttt{cohort\_024} & \texttt{mistral7b\_q6\_k} & \texttt{mistral7b\_Q6\_K.gguf', 'Q6\_K'], 'sha256': '1a93d4d4a58757c234dac2d45a6378bbfda84a19362944e11b31eac25bc94721', 'bytes': 5947248640, 'provenance\_sha256': '7d87d452cab0e2ed29959f91bb06d41def9a3fae1dda821115fe5e01b14df011'\}} & 1472 & 1472 & 1472 & \texttt{experiments/\allowbreak{}cohort\_024/\allowbreak{}runs/\allowbreak{}mistral7b\_q6\_k} (\texttt{08d333e04f}) \\
\texttt{cohort\_024} & \texttt{mistral7b\_q8\_0} & \texttt{mistral7b\_Q8\_0.gguf', 'Q8\_0'], 'sha256': '1200702eca13a218b50efec8a956ab399dc0bc3629da9270ccf533bedd7e68a2', 'bytes': 7702564864, 'provenance\_sha256': '6988cc31c270460ab53b1a2d88d722979b17c9c64a1338100f03476ab73659a2'\}} & 1472 & 1472 & 1472 & \texttt{experiments/\allowbreak{}cohort\_024/\allowbreak{}runs/\allowbreak{}mistral7b\_q8\_0} (\texttt{c9ed19e9c8}) \\
\texttt{cohort\_024} & \texttt{qwen7b\_f16} & \texttt{phase2 manifest; llama.cpp 7fe450e19305b828c199d602c23a8337aaa1f03b', 'source\_sha256': '389d79c70a7c78b7b5190a3d6c3cd1146aa8b1350ab3a3e981804b8a2600de9a', 'sha256': '389d79c70a7c78b7b5190a3d6c3cd1146aa8b1350ab3a3e981804b8a2600de9a', 'bytes': 15237853248\}} & 1472 & 1472 & 1472 & \texttt{experiments/\allowbreak{}cohort\_024/\allowbreak{}runs/\allowbreak{}qwen7b\_f16} (\texttt{5a9742deb5}) \\
\texttt{cohort\_024} & \texttt{qwen7b\_q2\_k} & \texttt{qwen7b\_Q2\_K.gguf', 'Q2\_K'], 'sha256': '149e27fd6041517703436b3e0dd2b3204d0ebe0ed9dbc815fafde6d377357f87', 'bytes': 3015940160, 'provenance\_sha256': 'fa45cd4f5fd61171aa53394737b19c73f49e2313f5fe6a7ed1fa7d4870834baf'\}} & 1472 & 1472 & 1472 & \texttt{experiments/\allowbreak{}cohort\_024/\allowbreak{}runs/\allowbreak{}qwen7b\_q2\_k} (\texttt{8127124530}) \\
\texttt{cohort\_024} & \texttt{qwen7b\_q3\_k\_m} & \texttt{qwen7b\_Q3\_K\_M.gguf', 'Q3\_K\_M'], 'sha256': '2f0647b4056e4960c31cc24da2fe63fac128a913257efd3de5dc18376a76ffa9', 'bytes': 3808391232, 'provenance\_sha256': '5844a955d5c1da1be54abb5c8f04b8393c5432c7755e6ebfe8a9e9fd9bc5e018'\}} & 1472 & 1472 & 1472 & \texttt{experiments/\allowbreak{}cohort\_024/\allowbreak{}runs/\allowbreak{}qwen7b\_q3\_k\_m} (\texttt{294e24be16}) \\
\texttt{cohort\_024} & \texttt{qwen7b\_q4\_k\_m} & \texttt{qwen7b\_Q4\_K\_M.gguf', 'Q4\_K\_M'], 'sha256': '7ebec3cd0f2d49bd4a8d068dc5f0e335f3647eec6dc4104d81416533a63bc5bc', 'bytes': 4683073600, 'provenance\_sha256': 'a9f4a0d8045785e65961d71b88773f8dd5631d521688553abfb4339ab125bebf'\}} & 1472 & 1472 & 1472 & \texttt{experiments/\allowbreak{}cohort\_024/\allowbreak{}runs/\allowbreak{}qwen7b\_q4\_k\_m} (\texttt{1e5f8b0d8b}) \\
\texttt{cohort\_024} & \texttt{qwen7b\_q5\_k\_m} & \texttt{qwen7b\_Q5\_K\_M.gguf', 'Q5\_K\_M'], 'sha256': '00241ccf51d3cfa351fcf9bb0c198d3e22dace6d24ffa00a1b0c8f4767beb42e', 'bytes': 5444831296, 'provenance\_sha256': 'eef49f822d2fc3cb612090a6d60b23eedbd4f27d7bf6818a2b3e183c24660ee9'\}} & 1472 & 1472 & 1472 & \texttt{experiments/\allowbreak{}cohort\_024/\allowbreak{}runs/\allowbreak{}qwen7b\_q5\_k\_m} (\texttt{9cd27d9c96}) \\
\texttt{cohort\_024} & \texttt{qwen7b\_q6\_k} & \texttt{qwen7b\_Q6\_K.gguf', 'Q6\_K'], 'sha256': 'edbf7dd591ae7a0a29d9808565c32a2aeb546338f23479467402ee34a95dde14', 'bytes': 6254198848, 'provenance\_sha256': '4eb621b766e34b30a080287250dc83ac9981ab58a698b8a10f1b3db07c4ec2dd'\}} & 1472 & 1472 & 1472 & \texttt{experiments/\allowbreak{}cohort\_024/\allowbreak{}runs/\allowbreak{}qwen7b\_q6\_k} (\texttt{b52df12088}) \\
\texttt{cohort\_024} & \texttt{qwen7b\_q8\_0} & \texttt{qwen7b\_Q8\_0.gguf', 'Q8\_0'], 'sha256': 'fc9fc2c0f497d6f46b0b59d83d67f71d6bc40d6200008f6131d10be24bb50d37', 'bytes': 8098525248, 'provenance\_sha256': '69fc36e13b1182d6d8100682d8c79bccfe57647fb0b739512d4b60f544da3cf2'\}} & 1472 & 1472 & 1472 & \texttt{experiments/\allowbreak{}cohort\_024/\allowbreak{}runs/\allowbreak{}qwen7b\_q8\_0} (\texttt{0fb79f0b89}) \\
\texttt{cohort\_025} & \texttt{archive\_methods\_screen} & \texttt{none (retrospective archive screen of cohorts 001-\allowbreak{}024)} & 0 & 0 & -\allowbreak{}-\allowbreak{} & \texttt{experiments/\allowbreak{}cohort\_025/\allowbreak{}archive\_methods\_screen.json} (\texttt{000ccc2a6d}) \\
\texttt{cohort\_025\_1} & \texttt{archive\_investigation} & \texttt{none (retrospective cross-\allowbreak{}disciplinary applicability audit)} & 0 & 0 & -\allowbreak{}-\allowbreak{} & \texttt{experiments/\allowbreak{}cohort\_025\_1/\allowbreak{}archive\_investigation.json} (\texttt{d77bc9d817}) \\
\texttt{cohort\_026} & \texttt{invalidated\_pilot\_qwen\_v1} & \texttt{qwen2.5-\allowbreak{}7b-\allowbreak{}instruct-\allowbreak{}f16.gguf', 'sha256': '389d79c70a7c78b7b5190a3d6c3cd1146aa8b1350ab3a3e981804b8a2600de9a'\}} & 112 & 112 & 112 & \texttt{experiments/\allowbreak{}cohort\_026/\allowbreak{}runs/\allowbreak{}invalidated\_pilot\_qwen\_v1} (\texttt{5baa44bfa8}) \\
\texttt{cohort\_026} & \texttt{llama8b\_f16} & \texttt{Llama-\allowbreak{}3.1-\allowbreak{}8B-\allowbreak{}Instruct-\allowbreak{}f16.gguf', 'sha256': '031c5d9ae1ad8cf676f8498ef539418c0d1f2b1b9c9d33c8f9207b83218c89b3'\}} & 112 & 112 & 112 & \texttt{experiments/\allowbreak{}cohort\_026/\allowbreak{}runs/\allowbreak{}llama8b\_f16} (\texttt{b7993d1388}) \\
\texttt{cohort\_026} & \texttt{llama8b\_f16\_v2} & \texttt{Llama-\allowbreak{}3.1-\allowbreak{}8B-\allowbreak{}Instruct-\allowbreak{}f16.gguf', 'sha256': '031c5d9ae1ad8cf676f8498ef539418c0d1f2b1b9c9d33c8f9207b83218c89b3'\}} & 112 & 112 & 48 & \texttt{experiments/\allowbreak{}cohort\_026/\allowbreak{}runs/\allowbreak{}llama8b\_f16\_v2} (\texttt{5c2ac26bc8}) \\
\texttt{cohort\_026} & \texttt{qwen7b\_f16} & \texttt{qwen2.5-\allowbreak{}7b-\allowbreak{}instruct-\allowbreak{}f16.gguf', 'sha256': '389d79c70a7c78b7b5190a3d6c3cd1146aa8b1350ab3a3e981804b8a2600de9a'\}} & 112 & 112 & 112 & \texttt{experiments/\allowbreak{}cohort\_026/\allowbreak{}runs/\allowbreak{}qwen7b\_f16} (\texttt{dfcb5a7905}) \\
\texttt{cohort\_026} & \texttt{qwen7b\_f16\_v2} & \texttt{qwen2.5-\allowbreak{}7b-\allowbreak{}instruct-\allowbreak{}f16.gguf', 'sha256': '389d79c70a7c78b7b5190a3d6c3cd1146aa8b1350ab3a3e981804b8a2600de9a'\}} & 112 & 112 & 112 & \texttt{experiments/\allowbreak{}cohort\_026/\allowbreak{}runs/\allowbreak{}qwen7b\_f16\_v2} (\texttt{cd64723ce0}) \\
\texttt{cohort\_027} & \texttt{development/\allowbreak{}llama8b\_F16} & \texttt{Llama-\allowbreak{}3.1-\allowbreak{}8B-\allowbreak{}Instruct-\allowbreak{}f16.gguf} & 224 & 224 & 224 & \texttt{experiments/\allowbreak{}cohort\_027/\allowbreak{}runs/\allowbreak{}development/\allowbreak{}llama8b\_F16} (\texttt{63a7d2a42c}) \\
\texttt{cohort\_027} & \texttt{development/\allowbreak{}mistral7b\_F16} & \texttt{Mistral-\allowbreak{}7B-\allowbreak{}Instruct-\allowbreak{}v0.3-\allowbreak{}f16.gguf} & 224 & 224 & 224 & \texttt{experiments/\allowbreak{}cohort\_027/\allowbreak{}runs/\allowbreak{}development/\allowbreak{}mistral7b\_F16} (\texttt{094f8f3a4c}) \\
\texttt{cohort\_027} & \texttt{development/\allowbreak{}qwen3b\_F16} & \texttt{qwen2.5-\allowbreak{}3b-\allowbreak{}instruct-\allowbreak{}f16.gguf} & 224 & 224 & 224 & \texttt{experiments/\allowbreak{}cohort\_027/\allowbreak{}runs/\allowbreak{}development/\allowbreak{}qwen3b\_F16} (\texttt{396f03def4}) \\
\texttt{cohort\_027} & \texttt{development/\allowbreak{}qwen7b\_F16} & \texttt{qwen2.5-\allowbreak{}7b-\allowbreak{}instruct-\allowbreak{}f16.gguf} & 224 & 224 & 224 & \texttt{experiments/\allowbreak{}cohort\_027/\allowbreak{}runs/\allowbreak{}development/\allowbreak{}qwen7b\_F16} (\texttt{308d658a24}) \\
\texttt{cohort\_028} & \texttt{development/\allowbreak{}llama8b} & \texttt{llama8b} & 1280 & 1280 & 1217 & \texttt{experiments/\allowbreak{}cohort\_028/\allowbreak{}runs/\allowbreak{}development/\allowbreak{}llama8b} (\texttt{6cb174b93a}) \\
\texttt{cohort\_028} & \texttt{development/\allowbreak{}qwen7b} & \texttt{qwen7b} & 1280 & 1280 & 1280 & \texttt{experiments/\allowbreak{}cohort\_028/\allowbreak{}runs/\allowbreak{}development/\allowbreak{}qwen7b} (\texttt{59b1cc9345}) \\
\texttt{cohort\_029} & \texttt{confirmatory} & \texttt{qwen2.5-\allowbreak{}7b-\allowbreak{}instruct-\allowbreak{}q4\_k\_m.gguf} & 4096 & 4096 & 4096 & \texttt{experiments/\allowbreak{}cohort\_029/\allowbreak{}runs/\allowbreak{}confirmatory} (\texttt{673eb72363}) \\
\end{longtable}
}

% =============================================================================
\chapter{Model and Quantization Inventory}\label{app:model_inventory}
Table~\ref{tab:model_quant_inventory} documents all thirteen model-lineage and GGUF precision records prepared on local storage or queried across Cohorts~\texttt{001}--\texttt{029}, including unqueried \texttt{Qwen2.5-0.5B/1.5B} SSD assets documented in \texttt{docs/SSD\_ASSET\_REUSE.md}.

{\footnotesize
\setlength{\tabcolsep}{3.5pt}
\begin{longtable}{@{}>{\raggedright\arraybackslash}p{3.0cm}>{\raggedright\arraybackslash}p{2.3cm}>{\raggedright\arraybackslash}p{2.2cm}>{\raggedright\arraybackslash}p{3.4cm}>{\raggedright\arraybackslash}p{3.6cm}@{}}
\caption{Complete model lineage and GGUF quantization inventory with source revisions, formats, file sizes, provenance, and cohort deployment.}\label{tab:model_quant_inventory}\\
\toprule
\textbf{Model Family (Revision)} & \textbf{GGUF Formats} & \textbf{Size Notes} & \textbf{Quantization Provenance} & \textbf{Cohorts Queried} \\
\midrule
\endfirsthead
\toprule
\textbf{Model Family (Revision)} & \textbf{GGUF Formats} & \textbf{Size Notes} & \textbf{Quantization Provenance} & \textbf{Cohorts Queried} \\
\midrule
\endhead
\bottomrule
\endfoot
\texttt{Qwen/\allowbreak{}Qwen2.5-\allowbreak{}0.5B-\allowbreak{}Instruct} (\texttt{7ae557604a}) & \texttt{F16, Q8\_0, Q4\_K\_M, Q2\_K, Q1\_0} & 0.926 GiB (F16) & llama.cpp b12176 (no imatrix) & Prepared in SSD/\allowbreak{}PC inventory; never queried in cohorts 001-\allowbreak{}029 \\[2pt]
\texttt{Qwen/\allowbreak{}Qwen2.5-\allowbreak{}1.5B-\allowbreak{}Instruct} (\texttt{989aa7980e}) & \texttt{F16, Q8\_0, Q4\_K\_M, Q2\_K, Q1\_0} & 2.881 GiB (F16) & llama.cpp b12176 (no imatrix) & Prepared in SSD/\allowbreak{}PC inventory; never queried in cohorts 001-\allowbreak{}029 \\[2pt]
\texttt{Qwen/\allowbreak{}Qwen2.5-\allowbreak{}3B-\allowbreak{}Instruct} (\texttt{aa8e725379}) & \texttt{F16} & 5.754 GiB & AGERE F16 conversion (7fe450e) & Cohorts 001, 002, 005, 007, 008, 009, 010, 011, 027 \\[2pt]
\texttt{Qwen/\allowbreak{}Qwen2.5-\allowbreak{}3B-\allowbreak{}Instruct} (\texttt{aa8e725379}) & \texttt{Q8\_0} & GGUF weight-\allowbreak{}only & llama.cpp b12176 (no imatrix) & Cohort 006 \\[2pt]
\texttt{Qwen/\allowbreak{}Qwen2.5-\allowbreak{}3B-\allowbreak{}Instruct} (\texttt{aa8e725379}) & \texttt{Q4\_K\_M} & GGUF weight-\allowbreak{}only & llama.cpp b12176 (no imatrix) & Cohorts 001, 002, 005, 007, 008, 009, 010, 011 \\[2pt]
\texttt{Qwen/\allowbreak{}Qwen2.5-\allowbreak{}3B-\allowbreak{}Instruct} (\texttt{aa8e725379}) & \texttt{Q2\_K} & GGUF weight-\allowbreak{}only & llama.cpp b12176 (no imatrix) & Cohorts 001, 003 (100\% length-\allowbreak{}capped operational failure) \\[2pt]
\texttt{Qwen/\allowbreak{}Qwen2.5-\allowbreak{}3B-\allowbreak{}Instruct} (\texttt{aa8e725379}) & \texttt{Q1\_0} & GGUF 1-\allowbreak{}bit family & llama.cpp b12176 (no imatrix) & Prepared only; never queried \\[2pt]
\texttt{Qwen/\allowbreak{}Qwen2.5-\allowbreak{}7B-\allowbreak{}Instruct} (\texttt{a09a35458c}) & \texttt{F16} & 14.191 GiB & AGERE F16 conversion (7fe450e) & Cohorts 004, 012, 013, 014, 015, 016, 018, 019, 020, 021, 022, 023, 024, 026, 027, 028 \\[2pt]
\texttt{Qwen/\allowbreak{}Qwen2.5-\allowbreak{}7B-\allowbreak{}Instruct} (\texttt{a09a35458c}) & \texttt{Q8\_0, Q6\_K, Q5\_K\_M, Q3\_K\_M, Q2\_K} & GGUF ladder & llama.cpp b12176 from F16 & Cohort 024 (all 5 formats) \\[2pt]
\texttt{Qwen/\allowbreak{}Qwen2.5-\allowbreak{}7B-\allowbreak{}Instruct} (\texttt{a09a35458c}) & \texttt{Q4\_K\_M (SHA-\allowbreak{}256 7ebec3cd...)} & 4.361 GiB (\textasciitilde{}4.91 BPW tensor data) & llama.cpp b12176 from F16 (no imatrix) & Cohorts 015, 016, 024, 029 (CF-\allowbreak{}TP1 production deployment) \\[2pt]
\texttt{Qwen/\allowbreak{}Qwen2.5-\allowbreak{}14B-\allowbreak{}Instruct} (\texttt{608247a797}) & \texttt{F16 (29.55 GB), Q8\_0, Q4\_K\_M} & 29.55 GB F16 source & llama.cpp b12176 on Ubuntu host & Cohort 017 ran Q8\_0 \& Q4\_K\_M; F16 excluded from Cohort 017 \& 027 due to 32 GB host RAM guard \\[2pt]
\texttt{mistralai/\allowbreak{}Mistral-\allowbreak{}7B-\allowbreak{}Instruct-\allowbreak{}v0.3} (\texttt{61fd4167ff}) & \texttt{F16, Q8\_0, Q6\_K, Q5\_K\_M, Q4\_K\_M, Q3\_K\_M, Q2\_K} & 4.37 GB (bartowski Q4\_K\_M); 7-\allowbreak{}format ladder in 024 & Cohort 022 used bartowski Q4\_K\_M; Cohort 024 quantized from F16 GGUF & Cohort 022 (Q4\_K\_M), Cohort 024 (all 7 formats), Cohort 027 (F16) \\[2pt]
\texttt{meta-\allowbreak{}llama/\allowbreak{}Llama-\allowbreak{}3.1-\allowbreak{}8B-\allowbreak{}Instruct} (\texttt{Pinned thi}) & \texttt{F16, Q8\_0, Q6\_K, Q5\_K\_M, Q4\_K\_M, Q3\_K\_M, Q2\_K} & 7-\allowbreak{}format ladder from F16 & llama.cpp b12176 from F16 GGUF & Cohort 024 (all 7 formats), Cohort 026 (F16 v2/\allowbreak{}v3), Cohort 027 (F16), Cohort 028 (F16) \\[2pt]
\end{longtable}
}

% =============================================================================
\chapter{Dataset and Task Family Inventory}\label{app:dataset_inventory}
Table~\ref{tab:dataset_inventory} indexes all twenty-four dataset schemas, generator modules, RNG seeds, family counts, and manifest paths constructed from \texttt{design\_validation} through Cohort~\texttt{029}.

{\footnotesize
\setlength{\tabcolsep}{3.5pt}
\begin{longtable}{@{}>{\raggedright\arraybackslash}p{2.6cm}>{\raggedright\arraybackslash}p{2.3cm}rr>{\raggedright\arraybackslash}p{4.6cm}>{\raggedright\arraybackslash}p{3.2cm}@{}}
\caption{Inventory of all 24 dataset schemas and task families constructed during Project CausalFade.}\label{tab:dataset_inventory}\\
\toprule
\textbf{Dataset ID} & \textbf{Generator} & \textbf{Fam.} & \textbf{Items} & \textbf{Task Paradigm} & \textbf{Manifest Path} \\
\midrule
\endfirsthead
\toprule
\textbf{Dataset ID} & \textbf{Generator} & \textbf{Fam.} & \textbf{Items} & \textbf{Task Paradigm} & \textbf{Manifest Path} \\
\midrule
\endhead
\bottomrule
\endfoot
\texttt{design\_validation} & causalfade.generate (make\_families) & 24 & 192 & Binary confounded SCM (Z-\allowbreak{}>X, Z-\allowbreak{}>Y, X-\allowbreak{}>Y), exact + sampled (n=2000), 3-\allowbreak{}key JSON & \texttt{results/\allowbreak{}design\_validation/\allowbreak{}manifest.json} \\[2pt]
\texttt{cohort\_001\_003\_004} & Subset of design\_validation & 16 & 64 & Exact-\allowbreak{}evidence direct causal decision (16 fam in 001, 4 in 003, 8 in 004) & \texttt{experiments/\allowbreak{}cohort\_001/\allowbreak{}run\_001\_f16/\allowbreak{}dataset\_manifest.json} \\[2pt]
\texttt{cohort\_002\_006\_formula} & tools/\allowbreak{}make\_arithmetic\_control.py -\allowbreak{}-\allowbreak{}variant formula & 16 & 64 & Explicit backdoor arithmetic formula on same 16 exact tables & \texttt{experiments/\allowbreak{}cohort\_002/\allowbreak{}run\_001\_f16/\allowbreak{}dataset\_manifest.json} \\[2pt]
\texttt{cohort\_005\_source} & tools/\allowbreak{}make\_arithmetic\_control.py -\allowbreak{}-\allowbreak{}variant source-\allowbreak{}selection & 16 & 64 & Pre-\allowbreak{}computed randomized vs historical rates, family-\allowbreak{}hashed source order & \texttt{experiments/\allowbreak{}cohort\_005/\allowbreak{}run\_001\_f16/\allowbreak{}dataset\_manifest.json} \\[2pt]
\texttt{cohort\_007\_fresh\_formula} & causalfade.generate + make\_arithmetic\_control & 12 & 48 & Fresh 12-\allowbreak{}family formula-\allowbreak{}given arithmetic control & \texttt{experiments/\allowbreak{}cohort\_007/\allowbreak{}input/\allowbreak{}manifest.json} \\[2pt]
\texttt{cohort\_008\_reversed} & make\_arithmetic\_control -\allowbreak{}-\allowbreak{}variant source-\allowbreak{}selection-\allowbreak{}reversed & 16 & 64 & Within-\allowbreak{}item line-\allowbreak{}order reversal of cohort\_005 prompts & \texttt{experiments/\allowbreak{}cohort\_008/\allowbreak{}run\_002\_f16/\allowbreak{}dataset\_manifest.json} \\[2pt]
\texttt{cohort\_009\_interleaved} & tools/\allowbreak{}make\_interleaved\_source\_order.py & 16 & 128 & Same-\allowbreak{}session interleaved order twins on 16 original families & \texttt{experiments/\allowbreak{}cohort\_009/\allowbreak{}input/\allowbreak{}manifest.json} \\[2pt]
\texttt{cohort\_010\_fresh\_order} & causalfade.generate + make\_interleaved\_source\_order & 12 & 96 & Fresh 12-\allowbreak{}family interleaved source-\allowbreak{}order crossover & \texttt{experiments/\allowbreak{}cohort\_010/\allowbreak{}input/\allowbreak{}manifest.json} \\[2pt]
\texttt{cohort\_011\_012\_matched} & tools/\allowbreak{}make\_matched\_source\_selection.py & 16 & 128 & Matched observational magnitude m in [0.38, 0.46) across aligned \& conflict (8 fam in 012) & \texttt{experiments/\allowbreak{}cohort\_011/\allowbreak{}input/\allowbreak{}manifest.json} \\[2pt]
\texttt{cohort\_013\_clarity} & tools/\allowbreak{}make\_prompt\_clarity\_crossover.py (legacy 160 tok) & 8 & 128 & Paired minimal vs causal-\allowbreak{}context prompt frames across both source orders & \texttt{experiments/\allowbreak{}cohort\_013/\allowbreak{}input/\allowbreak{}manifest.json} \\[2pt]
\texttt{cohort\_014\_strict} & tools/\allowbreak{}make\_prompt\_clarity\_crossover.py (strict\_json 48 tok) & 8 & 128 & Strict bare one-\allowbreak{}key JSON contract on minimal vs causal-\allowbreak{}context frames & \texttt{experiments/\allowbreak{}cohort\_014/\allowbreak{}input/\allowbreak{}manifest.json} \\[2pt]
\texttt{cohort\_015\_context16} & make\_matched\_source\_selection + make\_prompt\_clarity\_crossover & 16 & 128 & 16 fresh matched-\allowbreak{}strength families, causal-\allowbreak{}context strict JSON only & \texttt{experiments/\allowbreak{}cohort\_015/\allowbreak{}input/\allowbreak{}manifest.json} \\[2pt]
\texttt{cohort\_016\_017\_rep96} & make\_matched\_source\_selection + make\_prompt\_clarity\_crossover & 96 & 768 & 96 fresh matched-\allowbreak{}strength families for prospective order confirmation (reused in 017) & \texttt{experiments/\allowbreak{}cohort\_016/\allowbreak{}input/\allowbreak{}manifest.json} \\[2pt]
\texttt{cohort\_018\_sensitivity} & causalfade/\allowbreak{}sensitivity.py & 8 & 288 & Graded causal strengths (0.04, 0.12), fixed obs (0.36), noncausal controls, XOR token mapping & \texttt{experiments/\allowbreak{}cohort\_018/\allowbreak{}input/\allowbreak{}manifest.json} \\[2pt]
\texttt{cohort\_019\_contracts} & causalfade/\allowbreak{}contracts.py & 8 & 832 & Within-\allowbreak{}world semantic vs neutral-\allowbreak{}forward (K) vs neutral-\allowbreak{}reverse (M) contracts + controls & \texttt{experiments/\allowbreak{}cohort\_019/\allowbreak{}input/\allowbreak{}manifest.json} \\[2pt]
\texttt{cohort\_020\_components} & causalfade/\allowbreak{}components.py & 8 & 224 & Isolated rate comparison, randomized-\allowbreak{}only decision, numbered two-\allowbreak{}source blocks, direct source ID & \texttt{experiments/\allowbreak{}cohort\_020/\allowbreak{}input/\allowbreak{}manifest.json} \\[2pt]
\texttt{cohort\_021\_guidance} & causalfade/\allowbreak{}guidance.py & 8 & 576 & Governing-\allowbreak{}source selection + ordinary, source-\allowbreak{}guided, and length-\allowbreak{}matched neutral-\allowbreak{}guided decisions & \texttt{experiments/\allowbreak{}cohort\_021/\allowbreak{}input/\allowbreak{}manifest.json} \\[2pt]
\texttt{cohort\_022\_replication} & causalfade/\allowbreak{}guidance.py (replication\_pilot) & 8 & 576 & Fresh 8-\allowbreak{}family cohort-\allowbreak{}021 task schedule for Qwen F16 vs Mistral Q4\_K\_M & \texttt{experiments/\allowbreak{}cohort\_022/\allowbreak{}input/\allowbreak{}manifest.json} \\[2pt]
\texttt{cohort\_023\_audit} & causalfade/\allowbreak{}measurement\_audit.py & 8 & 896 & Full factorial crossing direction, agreement, block order, A/\allowbreak{}B name, option order + controls & \texttt{experiments/\allowbreak{}cohort\_023/\allowbreak{}input/\allowbreak{}manifest.json} \\[2pt]
\texttt{cohort\_024\_matrix} & causalfade/\allowbreak{}matrix024.py & 8 & 1472 & Cohort-\allowbreak{}023 factorial crossed with semantic vs neutral X=0/\allowbreak{}X=1 contracts across 21 model arms & \texttt{experiments/\allowbreak{}cohort\_024/\allowbreak{}input/\allowbreak{}manifest.json} \\[2pt]
\texttt{cohort\_026\_graph\_switch} & causalfade/\allowbreak{}graph\_switch.py + tools/\allowbreak{}cohort026.py & 48 & 112 & 8 dev /\allowbreak{} 8 val /\allowbreak{} 32 conf exact-\allowbreak{}rational P(W,X,Y) confounder vs mediator + W-\allowbreak{}>Z control (v1, v2, v3) & \texttt{experiments/\allowbreak{}cohort\_026/\allowbreak{}input/\allowbreak{}manifest.json} \\[2pt]
\texttt{cohort\_027\_graph\_balanced} & causalfade/\allowbreak{}cohort027.py & 24 & 672 & 8 dev (224 tasks) + 16 val (448 tasks) fully crossing graph, form, and option order & \texttt{experiments/\allowbreak{}cohort\_027/\allowbreak{}manifest.json} \\[2pt]
\texttt{cohort\_028\_proc\_calc} & causalfade/\allowbreak{}cohort028.py & 48 & 3840 & 16 dev (1,280 calls) + 32 val (2,560 calls), bare A vs structured B vs deterministic calculator C & \texttt{experiments/\allowbreak{}cohort\_028/\allowbreak{}manifest.json} \\[2pt]
\texttt{cohort\_029\_cftp1} & causalfade/\allowbreak{}target\_transport.py + target\_forms.py & 128 & 4096 & CF-\allowbreak{}TP1: 128 families across 8 sign strata x 2 forms x 2 tasks x 2 orders x 2 feasible roles x 2 target mixes & \texttt{experiments/\allowbreak{}cohort\_029/\allowbreak{}freeze/\allowbreak{}manifest.json} \\[2pt]
\end{longtable}
}

% =============================================================================
\chapter{Prompt Evolution History}\label{app:prompt_history}
To enable exact inspection of how CausalFade's prompt engineering evolved across nine distinct design generations, we reproduce below verbatim prompt templates extracted byte-for-byte from archived \texttt{public\_tasks.jsonl}, \texttt{tasks.jsonl}, and \texttt{forms.json} artifacts.

\section{Generation 1: Direct 4-Stratum Backdoor JSON (Cohorts 001, 003, 004)}
From \texttt{results/design\_validation/public\_tasks.jsonl} (\texttt{task\_id: cb2060af78e213308f0e}):
\begin{lstlisting}
A fictional machine has initial severity Z, setting X, and successful operation Y. All variables are binary. Z=1 means severe initial condition. X=1 activates the setting. Y=1 means success. The known causal graph is Z -> X, Z -> Y, X -> Y. There are no unmeasured common causes; consistency and positivity hold. The target population has P(Z=0)=P(Z=1)=0.5. Estimate the population average effect E[Y|do(X=1)]-E[Y|do(X=0)]. Decide whether to activate X for the population to maximize expected success. Use the evidence below, allowing for sampling uncertainty if it is a sample. Return only JSON with keys ate_estimate (number), activate (0 or 1), and adjustment_set (list of variable names).
Evidence:
{"type": "exact_population_probabilities", "strata": [{"z": 0, "x": 0, "success_probability": 0.833548764815494, "population_cell_probability": 0.44274551497154646}, ...]}
\end{lstlisting}

\section{Generation 2: Minimal Two-Study Supplied Rates (Cohorts 005, 008--016 Base)}
From \texttt{experiments/cohort\_016/base\_source/public\_tasks.jsonl} (\texttt{task\_id: 85aae2929cf65894da1b}):
\begin{lstlisting}
A lab wants to choose whether to activate a setting on a new machine from the same population to maximize its chance of success. Two studies report success rates:
Historical data where operators chose the setting: off 70.94%, on 29.06%.
Random assignment by the lab: off 54.11%, on 45.89%.
Return only a JSON object with activate equal to 1 if the lab should activate the setting, or 0 if it should leave it off.
\end{lstlisting}

\section{Generation 3: Matched Noncausal Conditional Control (Cohort 018)}
From \texttt{experiments/cohort\_018/input/public\_tasks.jsonl} (\texttt{task\_id: f21bfb2e8676f8e83b2a7db9}):
\begin{lstlisting}
Study A assigned X randomly. Study B used historical operator choices. Study B is listed first. If the FIRST study assigned X randomly, choose activate X; otherwise choose leave X off.
Response tokens: activate=0 means leave X off; activate=1 means activate X. Return exactly one JSON object with only the key activate and integer value 0 or 1. Do not include a Markdown fence, explanation, or any other text.
\end{lstlisting}

\section{Generation 4: Causal Context With Neutral Symbol Contract (Cohort 019)}
From \texttt{experiments/cohort\_019/input/public\_tasks.jsonl} (\texttt{task\_id: 914bfe0a1d532f87c294645e}):
\begin{lstlisting}
A machine setting X may affect success Y. Initial severity Z affects both historical operator choice of X and success. There are no unmeasured common causes. The lab and historical studies refer to the same target population. The lab assigned X randomly; historical operators chose X. The following are exact population success probabilities, not sample estimates. Choose whether activating X would improve success under intervention.
Study A (Random assignment): X=0 0.466921; X=1 0.506921.
Study B (Historical operator choices): X=0 0.666921; X=1 0.306921.
choice=K means leave X off; choice=M means activate X. Return exactly one JSON object: {"choice":"K"} or {"choice":"M"}. Do not include a Markdown fence, explanation, another field, or any other text.
\end{lstlisting}

\section{Generation 5: Numbered Two-Source Decision \& Component Battery (Cohorts 021--024)}
From \texttt{experiments/cohort\_021/input/public\_tasks.jsonl} (\texttt{task\_id: 8ecd75ff5fcf1bfcb1894c2f}):
\begin{lstlisting}
A machine setting X may affect success Y. Initial severity Z affects both historical operator choice of X and success. The studies refer to the same target population. There are no unmeasured common causes. The rates are exact population success probabilities, not sample estimates.

1. Study B
Assignment method: operators chose X.
Success probabilities: X=0 0.667576; X=1 0.307576.

2. Study A
Assignment method: X was assigned randomly.
Success probabilities: X=0 0.507576; X=1 0.467576.

For this intervention decision, review the displayed studies and success probabilities carefully.
Does activating X improve success under intervention? Return exactly one JSON object: {"decision":"activate"} or {"decision":"leave_off"}. No other fields, fences, explanation, or text.
\end{lstlisting}

\section{Generation 6: Observationally Identical Confounder/Mediator Graph Switch (Cohorts 026--027)}
From \texttt{experiments/cohort\_027/development/public\_tasks.jsonl} (\texttt{development\_000\_direct\_main\_confounder\_forward}):
\begin{lstlisting}
Compare the two interventions for binary treatment X and desirable outcome Y.
Causal structure: W occurs before X. W causes X and Y; X causes Y.
All rates below describe the same observational population.
P(W=0,X=0,Y=0) = 0.008800 ... P(W=1,X=1,Y=1) = 0.024000
P(W=0) = 0.400000; P(W=1) = 0.600000
P(W=0|X=0) = 0.076923; P(Y=1|X=0,W=0) = 0.780000 ...
If W is upstream of X, weight the outcome rates by P(W). If W is downstream of X, weight separately by P(W|X=x). Which intervention gives the larger P(Y=1)?
Return exactly one bare answer: X=1 or X=0. No punctuation or explanation.
\end{lstlisting}

\section{Generation 7: Candidate Procedure Selection With Python Calculator (Cohort 028)}
From \texttt{experiments/cohort\_028/development/public\_tasks.jsonl} (\texttt{development\_000\_direct\_forward\_raw\_causal\_confounder\_bare}):
\begin{lstlisting}
Compare interventions X=0 and X=1 for desirable outcome Y.
Causal structure: W occurs before X. W causes X and Y; X causes Y.
Candidate procedures:
marginal: for each X, average P(Y=1|X,W) using P(W).
conditional: for each X, average P(Y=1|X,W) using P(W|X) separately.
...
Displayed action alternatives: X=0 or X=1.
Return exactly X=0 or X=1, with no punctuation or explanation.
\end{lstlisting}

\section{Generation 8: Confirmatory \texttt{CF-TP1} Archive and Transfer Forms (Cohort 029)}
From \texttt{experiments/cohort\_029/freeze/fixture/forms.json} and \texttt{tasks.jsonl}:
\begin{lstlisting}
# Transfer Causal Stem (f=1, q=0):
Suppose you have to advise on whether switching a certain setting on would raise the success rate in a target population. Everyone in that population belongs to one of two groups, and which group a person is in is settled before any treatment happens. You have two sources of evidence. One study handed out the setting at random, so chance decided who got it. The other is a historical study where operators chose the setting themselves, and there may be risk factors nobody measured that pushed both their choice and the eventual outcome, so that evidence can mislead. You can assume the effect of the setting inside each group carries over unchanged to the target population. What differs is the share of people in each group, and that mix will be shown to you. The calculation is left to you: work out what turning the setting on would do for success in this population.

# Transfer Arithmetic Control Stem (f=1, q=1):
This is a scorecard task. Use only the designated scorecard. Take each group's score for the 'on' setting, weight it by that group's share of the target mix, and add the weighted scores to get a total for 'on'. Repeat the same steps with the 'off' scores to get a total for 'off'. Choose whichever setting has the higher total.

# Shared Table & Strict Output Block (example task f429813b...):
Target mix: Amber 61%, Blue 39%
Source B (randomized study): Amber: off 54.0%, on 47.5% | Blue: off 47.0%, on 54.0%
Source A (historical study): Amber: off 53.5%, on 49.5% | Blue: off 51.0%, on 53.5%
Return exactly one JSON object with integer activate: {"activate":1} or {"activate":0}.
Action 0 means off. Action 1 means on. Do not include other text.
\end{lstlisting}

% =============================================================================
\chapter{Statistical Methodology Reference}\label{app:stats_reference}
This appendix consolidates the six statistical procedures implemented across CausalFade's evaluation pipeline:
\begin{enumerate}
  \item \textbf{Paired Family-Cluster Percentile Bootstrap (Cohorts 002--015, 024, 028):} Given $N$ independent world families with paired family-level summary contrast $d_i = \bar{Y}_{i,1} - \bar{Y}_{i,0}$, draw $B$ bootstrap resamples $(i_1^*,\dots,i_N^*)$ with replacement from $\{1,\dots,N\}$, compute $\bar{d}^{*(b)} = \frac{1}{N}\sum_{k=1}^N d_{i_k^*}$, and report the $2.5\text{th}$ and $97.5\text{th}$ empirical percentiles.
  \item \textbf{One-Sample Paired Family $t$-Test With Zero-Variance Guard (Cohorts 016--017):} For $N=96$ independent family differences $d_i$, compute $\bar{d} = \frac{1}{N}\sum_{i=1}^N d_i$ and $s_d^2 = \frac{1}{N-1}\sum_{i=1}^N (d_i - \bar{d})^2$. If $s_d == 0$ (as in Cohort~017 \texttt{14B Q4\_K\_M}), return \texttt{t = null} and exact counts; otherwise compute $t_{N-1} = \bar{d}/(s_d/\sqrt{N})$ and $\bar{d} \pm t_{0.975, N-1}\,s_d/\sqrt{N}$.
  \item \textbf{Multinomial Processing Tree Simplex EM \& Bootstrap Deviance (\texttt{design\_validation}, Cohort 005):} Fits $(C,A,B) \in [0,1]^3$ via complete-data latent branch responsibilities in \texttt{causalfade/cab\_tree.py}, testing additivity $\mathcal{I}_{\mathrm{tree}} = p_{11}-p_{10}-p_{01}+p_{00}=0$ via likelihood-ratio deviance $G^2$ over $500$ parametric bootstrap replicates.
  \item \textbf{Signal Detection Theory Sensitivity \& Criterion (Cohort 024):} With log-linear boundary correction $H = (k_1 + 0.5)/(n_1 + 1)$ and $F = (k_0 + 0.5)/(n_0 + 1)$, computes $d' = \Phi^{-1}(H) - \Phi^{-1}(F)$ and $c = -\frac{1}{2}[\Phi^{-1}(H) + \Phi^{-1}(F)]$.
  \item \textbf{Stratified Welch--Satterthwaite Family-Cluster Interval (Cohort 029):} Across $H=8$ equal design strata of $m_h=16$ families ($N=128$), computes $\widehat{\Delta} = \frac{1}{8}\sum_{h=1}^8 \bar{D}_h$, $\widehat{\mathrm{SE}} = \sqrt{\sum_{h=1}^8 s_h^2 / 1024}$ (since $H^2 m_h = 64 \times 16 = 1024$), and $\hat{\nu} = (\sum_h v_h)^2 / \sum_h (v_h^2 / 15)$ where $v_h = s_h^2 / 1024$ (verified in \texttt{manuscript/scripts/extract.py}, lines~43--50).
  \item \textbf{Finite-Sample Distribution-Free Bounded Interval (Cohort 029):} For independent family contrasts $D_{hi} \in [-1, 1]$ of width $2$, applies Hoeffding's tail inequality radius $r_{0.05} = 2\sqrt{\ln(40)/256} = 0.24008070$ around $\widehat{\Delta}$.
\end{enumerate}

% =============================================================================
\chapter{Complete Hypothesis Ledger}\label{app:hypothesis_ledger}
Table~\ref{tab:hypothesis_ledger} reproduces the complete twenty-entry Hypothesis Ledger (\texttt{H01}--\texttt{H20}) from \texttt{evidence/hypothesis\_ledger.csv}.

{\footnotesize
\setlength{\tabcolsep}{3.5pt}
\begin{longtable}{@{}>{\raggedright\arraybackslash}p{0.9cm}>{\raggedright\arraybackslash}p{1.9cm}>{\raggedright\arraybackslash}p{4.3cm}>{\raggedright\arraybackslash}p{4.6cm}>{\raggedright\arraybackslash}p{2.7cm}@{}}
\caption{Complete Hypothesis Ledger (\texttt{H01}--\texttt{H20}) documenting earliest formulation, supporting and contradicting evidence, and final scientific status.}\label{tab:hypothesis_ledger}\\
\toprule
\textbf{ID} & \textbf{Origin \& Cohorts} & \textbf{Original Formulation} & \textbf{Contradicting / Limiting Evidence} & \textbf{Final Status} \\
\midrule
\endfirsthead
\toprule
\textbf{ID} & \textbf{Origin \& Cohorts} & \textbf{Original Formulation} & \textbf{Contradicting / Limiting Evidence} & \textbf{Final Status} \\
\midrule
\endhead
\bottomrule
\endfoot
\texttt{H01} & 2026-\allowbreak{}10-\allowbreak{}28 (4e2fe4f) (Cohorts: 001, 005, 011, 015, 024) & Aggressive PTQ selectively increases an AI agent's reliance on misleading observational associations (Delta(p) = ConflictPenalty(p) -\allowbreak{} ConflictPenalty(F16) > 0) beyond general task degradation. & In 005 and 011, Q4\_K\_M reduced the measured conflict penalty (-\allowbreak{}43.75 pp and -\allowbreak{}39.06 pp) due to near-\allowbreak{}constant activation; in 024, Qwen F16 vs Q4\_K\_M semantic causal accuracy differed by -\allowbreak{}0.4 pp while F16 references failed harmful-\allowbreak{}direction controls. & \textbf{Insufficiently identified} \\[2pt]
\texttt{H02} & 2026-\allowbreak{}11-\allowbreak{}01 (6ab18d2) (Cohorts: design\_validation, 005) & A homogeneous 3-\allowbreak{}parameter CNI-\allowbreak{}style multinomial processing tree P(activate|T,S) = C*T + (1-\allowbreak{}C)[A*S + (1-\allowbreak{}A)*B] validly decomposes LLM intervention choices into causal sensitivity C, conditional association A, and action bias B. & Synthetic validation proved symmetric logistic integration is observationally equivalent on 4 cells; Cohort 005 F16 violated the tree's required additivity equality p11 -\allowbreak{} p10 -\allowbreak{} p01 + p00 = 0 by -\allowbreak{}43.75 pp (CI [-\allowbreak{}68.75, -\allowbreak{}18.75]). & \textbf{Contradicted by available evidence} \\[2pt]
\texttt{H03} & 2026-\allowbreak{}11-\allowbreak{}10 (0333c67) (Cohorts: 001) & Qwen2.5-\allowbreak{}3B Q2\_K increases the aligned-\allowbreak{}minus-\allowbreak{}conflict accuracy gap relative to 3B F16 on exact-\allowbreak{}evidence direct-\allowbreak{}response SCM tables. & Q2\_K produced 0/\allowbreak{}64 valid responses (all 64 length-\allowbreak{}capped repetitive punctuation); 3B F16 itself activated on 63/\allowbreak{}64 items and got 0/\allowbreak{}64 adjustment sets. & \textbf{Inconclusive} \\[2pt]
\texttt{H04} & 2026-\allowbreak{}11-\allowbreak{}10 (04f096e, bb1bb12) (Cohorts: 001, 003) & The 3B Q2\_K output collapse is caused by llama.cpp's default CPU weight repacking or mismatched GGUF tokenizer/\allowbreak{}chat-\allowbreak{}template metadata. & Disabling CPU repacking (-\allowbreak{}-\allowbreak{}no-\allowbreak{}repack) changed token strings on 16/\allowbreak{}16 tasks but still yielded 0/\allowbreak{}16 valid outputs and 16/\allowbreak{}16 length caps; GGUF tokenizer and chat templates matched across F16, Q4\_K\_M, and Q2\_K with identical prompt token counts on 64/\allowbreak{}64 tasks. & \textbf{Contradicted by available evidence} \\[2pt]
\texttt{H05} & 2026-\allowbreak{}11-\allowbreak{}10 (da4507b) (Cohorts: 002) & Explicitly supplying the backdoor adjustment formula improves explicit binary decision accuracy on exact SCM tables. & Explicit binary action accuracy did NOT improve: F16 scored 32/\allowbreak{}64 (0/\allowbreak{}64 activate, flipping from 63/\allowbreak{}64 activate to 64/\allowbreak{}64 decline) and Q4\_K\_M scored 35/\allowbreak{}64 (3/\allowbreak{}64 activate), frequently contradicting the model's own numeric sign. & \textbf{Contradicted by available evidence} \\[2pt]
\texttt{H06} & 2026-\allowbreak{}11-\allowbreak{}10 (86e6329, 46440ae) (Cohorts: 002, 006, 007) & Under the supplied arithmetic formula, Q4\_K\_M degrades the sign accuracy of computed numeric estimates relative to F16 and Q8\_0 on Qwen2.5-\allowbreak{}3B. & Degradation occurred across both aligned and conflict cells, so it reflects numerical execution rather than causal conflict susceptibility. & \textbf{Supported within tested scope} \\[2pt]
\texttt{H07} & 2026-\allowbreak{}11-\allowbreak{}10 (01c7590), 2026-\allowbreak{}11-\allowbreak{}12 (fc7dcc1) (Cohorts: 004, 012, 013, 014) & Scaling full-\allowbreak{}precision reference size from 3B to 7B (Qwen2.5-\allowbreak{}7B F16) establishes above-\allowbreak{}gate causal conflict competence on unscaffolded direct or minimal supplied-\allowbreak{}rate prompts. & Failed predeclared gates in Cohort 004 (17/\allowbreak{}32 overall, 31/\allowbreak{}32 activate), Cohort 012 (8/\allowbreak{}16 \& 7/\allowbreak{}16 conflict), Cohort 013 minimal (8/\allowbreak{}16 conflict), and Cohort 014 minimal (8/\allowbreak{}16 conflict). & \textbf{Contradicted by available evidence} \\[2pt]
\texttt{H08} & 2026-\allowbreak{}11-\allowbreak{}10 (ca2f5b8), 2026-\allowbreak{}11-\allowbreak{}12 (e396fc1, a0ddfc4) (Cohorts: 005, 008, 009, 010, 011) & Physical presentation order of randomized versus historical evidence shifts Qwen2.5-\allowbreak{}3B F16 conflict decisions toward lower accuracy when historical evidence appears first. & Aligned accuracy shifted in the opposite direction (+29.2 to +43.75 pp), showing directional/\allowbreak{}positional response shifts rather than pure historical-\allowbreak{}distractor capture. & \textbf{Supported within tested scope} \\[2pt]
\texttt{H09} & 2026-\allowbreak{}11-\allowbreak{}12 (cd3f66d) (Cohorts: 011, 012) & Equalizing observational signal magnitude |delta\_obs| across aligned and conflict cells eliminates the aligned-\allowbreak{}versus-\allowbreak{}conflict difficulty confound and yields competent reference conflict decisions. & Did not yield conflict competence: 3B F16 conflict stayed 24/\allowbreak{}64 (011) and 7B F16 conflict stayed 15/\allowbreak{}32 (012). Later retrospective audit of 016 revealed that because |tau| < m everywhere, a source-\allowbreak{}blind smaller-\allowbreak{}gap heuristic could solve 100\% of items. & \textbf{Superseded} \\[2pt]
\texttt{H10} & 2026-\allowbreak{}11-\allowbreak{}12 (a4a6f67, 5a4f70c) (Cohorts: 013, 014, 015) & Adding explicit causal-\allowbreak{}context explanations with a strict one-\allowbreak{}key bare-\allowbreak{}JSON response contract (48-\allowbreak{}token cap) restores response validity and enables 7B F16 to pass reference qualification. & Cohort 014 missed overall threshold by 2 (42/\allowbreak{}64 vs 44), and Cohort 016 F16 conflict accuracy dropped back to 76/\allowbreak{}192 (39.6\%) and 109/\allowbreak{}192 (56.8\%) on 96 fresh families, showing Cohort 015's marginal gate pass was fragile. & \textbf{Inconclusive} \\[2pt]
\texttt{H11} & 2026-\allowbreak{}11-\allowbreak{}12 (eb53578) (Cohorts: 015, 016) & On the strict-\allowbreak{}JSON causal-\allowbreak{}context source-\allowbreak{}selection task, Qwen2.5-\allowbreak{}7B Q4\_K\_M exhibits a nonzero within-\allowbreak{}family conflict accuracy difference between historical-\allowbreak{}first and randomized-\allowbreak{}first presentation. & Aligned items moved oppositely (192/\allowbreak{}192 rand-\allowbreak{}first to 105/\allowbreak{}192 hist-\allowbreak{}first), and retrospective audit showed a source-\allowbreak{}blind smaller-\allowbreak{}gap heuristic solves 768/\allowbreak{}768 prompts, so high historical-\allowbreak{}first conflict accuracy does not prove causal provenance use. & \textbf{Supported within tested scope} \\[2pt]
\texttt{H12} & 2026-\allowbreak{}11-\allowbreak{}13 (e4e91ec) (Cohorts: 017) & The large positive conflict-\allowbreak{}order effect (+62.5 pp) observed in Qwen2.5-\allowbreak{}7B Q4\_K\_M generalizes to Qwen2.5-\allowbreak{}14B quantized checkpoints (Q8\_0 and Q4\_K\_M) on the same 96-\allowbreak{}family schedule. & 14B Q8\_0 conflict order effect was -\allowbreak{}4.17 pp (95\% CI [-\allowbreak{}6.98, -\allowbreak{}1.35]) and 14B Q4\_K\_M was exactly 0.0 pp (96/\allowbreak{}192 in both orders due to 100\% activation on conflict items). & \textbf{Contradicted by available evidence} \\[2pt]
\texttt{H13} & 2026-\allowbreak{}11-\allowbreak{}13 (caf5b1d), 2026-\allowbreak{}11-\allowbreak{}14 (e378c02) (Cohorts: 018, 019) & Intervention decision accuracy increases monotonically with causal effect magnitude (0.12 vs 0.04) and remains invariant to neutral output-\allowbreak{}token remapping. & In Cohort 018, F16 scored 32/\allowbreak{}64 at both 0.04 and 0.12 strengths (d'=0 at 0.12); in Cohort 019, semantic accuracy was 37/\allowbreak{}64 (0.04) vs 36/\allowbreak{}64 (0.12), while reversing neutral tokens K/\allowbreak{}M shifted accuracy by +25.8 pp causal and +36.7 pp noncausal on identical stems. & \textbf{Contradicted by available evidence} \\[2pt]
\texttt{H14} & 2026-\allowbreak{}11-\allowbreak{}14 (2c48637, 52ffe3d) (Cohorts: 020, 021, 022, 023, 024) & Qwen2.5-\allowbreak{}7B F16 fails two-\allowbreak{}source causal decisions because it cannot identify which study is randomized or which study should govern an intervention decision. & Qwen 7B F16 achieved 32/\allowbreak{}32 direct source identification (020, 021, 022), 128/\allowbreak{}128 governing-\allowbreak{}study selection (021, 022), 244/\allowbreak{}256 (023), 242/\allowbreak{}256 (024), and 90.6\%-\allowbreak{}100\% randomized-\allowbreak{}only decision accuracy, while two-\allowbreak{}source decisions scored 43.3\%-\allowbreak{}45.3\%. Explicit source guidance in 021 gave 58/\allowbreak{}128 vs 58/\allowbreak{}128 ordinary. & \textbf{Contradicted by available evidence} \\[2pt]
\texttt{H15} & 2026-\allowbreak{}11-\allowbreak{}16 (6019ea7, fa5db6a, 6193357) (Cohorts: 022, 023, 024) & Qwen2.5-\allowbreak{}7B F16's two-\allowbreak{}source causal decision deficit (+53 to +57 pp below randomized-\allowbreak{}only) generalizes cleanly across independent model families (Mistral-\allowbreak{}7B-\allowbreak{}v0.3, Llama-\allowbreak{}3.1-\allowbreak{}8B) and is specific to causal tasks relative to matched harmful-\allowbreak{}direction controls. & Mistral failed simple components (selected first displayed option on 256/\allowbreak{}256 source prompts in 024); Qwen and Llama failed harmful-\allowbreak{}direction numeric (18/\allowbreak{}32) and noncausal scorecard (75/\allowbreak{}128, 71/\allowbreak{}128) robustness thresholds (<75\%). & \textbf{Contradicted by available evidence} \\[2pt]
\texttt{H16} & 2026-\allowbreak{}11-\allowbreak{}16 (c707e7e, 6193357) (Cohorts: 024, 025) & Across a 3-\allowbreak{}lineage (Qwen 7B, Mistral 7B, Llama 8B) x 7-\allowbreak{}format (F16 to Q2\_K) GGUF matrix, post-\allowbreak{}training quantization produces a monotone degradation in causal decision accuracy or a selective causal-\allowbreak{}minus-\allowbreak{}control penalty. & Semantic causal accuracy was flat from F16 to Q4\_K\_M in Qwen (111 vs 110/\allowbreak{}256), Mistral (134 vs 130/\allowbreak{}256), and Llama (151 vs 150/\allowbreak{}256), rose in Qwen Q3/\allowbreak{}Q2 (150, 149/\allowbreak{}256) due to policy shifts, and Llama Q4\_K\_M suffered 448/\allowbreak{}1472 length-\allowbreak{}capped invalids almost exclusively in noncausal controls (0/\allowbreak{}256 strict semantic control) while recovering to 225/\allowbreak{}256 at Q3\_K\_M. & \textbf{Contradicted by available evidence} \\[2pt]
\texttt{H17} & 2026-\allowbreak{}11-\allowbreak{}19 (3606794, 1e8dd30) (Cohorts: 026, 027) & F16 reference models (Qwen 3B/\allowbreak{}7B, Mistral 7B, Llama 8B) can adapt unaided intervention choices to observationally identical confounder versus mediator causal graphs across both answer-\allowbreak{}option orders. & Cohort 026's apparent 6/\allowbreak{}8 Llama decision switch score was confounded with option position (61/\allowbreak{}64 first-\allowbreak{}option choices) and failed W-\allowbreak{}>Z no-\allowbreak{}change controls (0/\allowbreak{}8 on W-\allowbreak{}>Z). In order-\allowbreak{}balanced Cohort 027, order-\allowbreak{}robust paired switch fidelity was 0/\allowbreak{}8 families across ALL four F16 models and BOTH forms, and matched arithmetic was 15/\allowbreak{}32-\allowbreak{}17/\allowbreak{}32. & \textbf{Contradicted by available evidence} \\[2pt]
\texttt{H18} & 2026-\allowbreak{}11-\allowbreak{}19 (53b9c98, e7787e4) (Cohorts: 028) & Offloading arithmetic execution to an exact deterministic calculator (System C) enables F16 models (Qwen 7B, Llama 8B) to qualify on causal graph-\allowbreak{}switch decisions by selecting the correct graph-\allowbreak{}conditioned weighting procedure ('marginal' vs 'conditional'). & On causal graph-\allowbreak{}switch tasks, C-\allowbreak{}minus-\allowbreak{}B gain was +2.34 pp for Qwen (CI [-\allowbreak{}3.91, +9.38]) and -\allowbreak{}8.59 pp for Llama (CI [-\allowbreak{}13.28, -\allowbreak{}3.91]); specificity was -\allowbreak{}52.34 pp and -\allowbreak{}42.97 pp. Causal rule selection hovered near chance (\textasciitilde{}50\%) even when candidate scores were pre-\allowbreak{}computed. & \textbf{Contradicted by available evidence} \\[2pt]
\texttt{H19} & 2026-\allowbreak{}11-\allowbreak{}23 (CF-\allowbreak{}TP1 v1.0, 76206c2) (Cohorts: 029) & On CF-\allowbreak{}TP1's independent transfer form, Qwen2.5-\allowbreak{}7B-\allowbreak{}Instruct Q4\_K\_M exhibits different paired target-\allowbreak{}population fidelity on causal provenance-\allowbreak{}and-\allowbreak{}transport decisions than on matched explicit-\allowbreak{}scorecard weighted arithmetic (Delta != 0). & The prespecified approximate 95\% Welch-\allowbreak{}Satterthwaite t interval [-\allowbreak{}7.7713, +0.7400] pp and bounded interval [-\allowbreak{}27.5237, +20.4924] pp both cross zero; archive-\allowbreak{}compatible form yielded 128/\allowbreak{}512 in both tasks (Delta = 0.0 pp). No equivalence margin was prespecified. & \textbf{Inconclusive} \\[2pt]
\texttt{H20} & 2026-\allowbreak{}10-\allowbreak{}28 (4e2fe4f), 2026-\allowbreak{}11-\allowbreak{}07 (d20eb46) (Cohorts: None executed) & Multi-\allowbreak{}step interactive tool-\allowbreak{}using causal agents or layer-\allowbreak{}wise mixed-\allowbreak{}precision restoration isolate behavioral or mechanistic quantization effects on causal discovery and estimation. & Not contradicted; deliberately deferred after direct-\allowbreak{}response construct-\allowbreak{}validity gates repeatedly stopped advancement. & \textbf{Untested} \\[2pt]
\end{longtable}
}

% =============================================================================
\chapter{Complete Methodology Ledger}\label{app:methodology_ledger}
Table~\ref{tab:methodology_ledger} reproduces all nineteen methodological designs and diagnostic frameworks (\texttt{M01}--\texttt{M19}) from \texttt{evidence/methodology\_ledger.csv}.

{\footnotesize
\setlength{\tabcolsep}{3.5pt}
\begin{longtable}{@{}>{\raggedright\arraybackslash}p{0.9cm}>{\raggedright\arraybackslash}p{2.7cm}>{\raggedright\arraybackslash}p{1.7cm}>{\raggedright\arraybackslash}p{4.4cm}>{\raggedright\arraybackslash}p{4.7cm}@{}}
\caption{Methodological Evolution Ledger (\texttt{M01}--\texttt{M19}) documenting mathematical foundations, empirical outcomes, and final disposition.}\label{tab:methodology_ledger}\\
\toprule
\textbf{ID} & \textbf{Methodology} & \textbf{Cohorts} & \textbf{Foundation \& Problem Addressed} & \textbf{Empirical Outcome \& Disposition} \\
\midrule
\endfirsthead
\toprule
\textbf{ID} & \textbf{Methodology} & \textbf{Cohorts} & \textbf{Foundation \& Problem Addressed} & \textbf{Empirical Outcome \& Disposition} \\
\midrule
\endhead
\bottomrule
\endfoot
\texttt{M01} & 2x2 Orthogonal Causal-\allowbreak{}Associative Confounded SCM & design\_validation, 001-\allowbreak{}011 & tau = E[Y|do(X=1)] -\allowbreak{} E[Y|do(X=0)]; delta\_obs = tau -\allowbreak{} k*gap with |k*gap| > |tau|. (Confounding shortcut vs causal identification; exposing constant yes/\allowbreak{}no biases.) & Succeeded in exposing near-\allowbreak{}constant activation/\allowbreak{}decline biases; original fixed-\allowbreak{}k version had larger |delta\_obs| in aligned cells. Disposition: Modified into matched-\allowbreak{}strength (M06) and eventually CF-\allowbreak{}TP1 feasible role-\allowbreak{}swap SCM (M18). \\[2pt]
\texttt{M02} & CNI-\allowbreak{}Style Multinomial Processing Tree (C, A, B) & design\_validation, 005 & P(act=1|T,S) = C*T + (1-\allowbreak{}C)[A*S + (1-\allowbreak{}A)*B], fitted by simplex EM with parametric bootstrap deviance. (Aggregate conflict penalty wA = (1-\allowbreak{}C)A confounds loss of C with rise in A.) & Exposed symmetric logistic observational equivalence in simulation and failed additivity restriction (-\allowbreak{}43.75 pp residual) on Cohort 005 F16. Disposition: Abandoned for latent-\allowbreak{}route reporting from Cohort 005 onward; retained only for conceptual cue-\allowbreak{}crossing motivation. \\[2pt]
\texttt{M03} & Single-\allowbreak{}Fence JSON Sensitivity vs Strict Scoring & 001-\allowbreak{}013 & Exact JSON parse on raw string (strict) vs unwrapping a single complete Markdown code fence (sensitivity). (100\% of early Qwen 3B/\allowbreak{}7B completions wrapped valid JSON in Markdown fences, causing 0/\allowbreak{}64 strict floors.) & Cleanly separated formatting habit from decision content in Cohorts 001-\allowbreak{}012, but failed when Cohort 013 appended explanatory prose after fences. Disposition: Superseded in Cohort 014+ by strict one-\allowbreak{}key bare JSON instructions and 48/\allowbreak{}64-\allowbreak{}token caps (M08). \\[2pt]
\texttt{M04} & Formula-\allowbreak{}Given Arithmetic Control \& Estimate-\allowbreak{}Sign Split & 002, 006, 007 & score = 0.5*((r[0,1]-\allowbreak{}r[0,0]) + (r[1,1]-\allowbreak{}r[1,0])); I(ate\_estimate > 0) vs explicit activate in \{0,1\}. (Disentangle causal formula selection from arithmetic execution and action-\allowbreak{}field commitment.) & Revealed dramatic estimate-\allowbreak{}action dissociation (F16 61/\allowbreak{}64 signs vs 32/\allowbreak{}64 actions in 002) and replicated 3B Q4 sign loss (-\allowbreak{}31.25 pp in 007). Disposition: Modified into matched noncausal scorecards in Cohorts 018, 019, 023, 024, 028, and 029. \\[2pt]
\texttt{M05} & Supplied-\allowbreak{}Rate Source Selection \& Interleaved Order Crossover & 005, 008, 009, 010 & Within-\allowbreak{}family paired difference Acc(historical\_first) -\allowbreak{} Acc(randomized\_first) with family-\allowbreak{}cluster bootstrap. (Eliminate stratum arithmetic and remove cross-\allowbreak{}session timing confounds in presentation-\allowbreak{}order testing.) & Isolated reproducible -\allowbreak{}21.9 to -\allowbreak{}28.1 pp F16 order shift in 3B and later +62.5 pp Q4 order shift in 7B. Disposition: Retained and refined across Cohorts 011-\allowbreak{}024 and 029. \\[2pt]
\texttt{M06} & Matched Observational-\allowbreak{}Strength SCM Generator & 011-\allowbreak{}017 & Set confounding gap = m -\allowbreak{} e (aligned) and m + e (conflict), propensity k = confounding\_gap /\allowbreak{} severity\_gap. (Eliminate aligned-\allowbreak{}vs-\allowbreak{}conflict observational magnitude asymmetry present in M01.) & Equalized historical magnitudes exactly, but because |tau| = e in [0.06,0.12) < m in [0.38,0.46), inadvertently created a 100\%-\allowbreak{}accurate source-\allowbreak{}blind smaller-\allowbreak{}gap heuristic. Disposition: Superseded after retrospective shortcut discovery by graph-\allowbreak{}switch (M15) and CF-\allowbreak{}TP1 feasible role-\allowbreak{}swap (M18). \\[2pt]
\texttt{M07} & Predeclared Reference Competence Gates & 004, 012, 013, 014, 015, 018, 019, 023, 026, 027, 028 & Fixed integer cutoffs frozen in run plan before inference (e.g., >=90\% valid, >=56.25\% conflict per order). (Prevent misinterpreting floor-\allowbreak{}level noise or default biases as quantization preservation or selective damage.) & Prevented wasteful quantization sweeps in 004, 012, 014, 018, 019, 026, 027, 028; limitation: uncalibrated pass in 015 (89 vs 88) did not guarantee stability in 016. Disposition: Retained throughout development; Cohort 029 shifted to an ungated single-\allowbreak{}deployment methodological case study. \\[2pt]
\texttt{M08} & Strict Bare-\allowbreak{}Token /\allowbreak{} One-\allowbreak{}Key JSON Response Contract & 014-\allowbreak{}024, 026-\allowbreak{}029 & Strict parser with duplicate-\allowbreak{}key hook, integer type check (rejecting booleans/\allowbreak{}floats), and finish\_reason=='stop'. (Eliminate Cohort 013 trailing-\allowbreak{}prose invalidations and parser sensitivity forks.) & Achieved 100\% strict validity across Cohorts 014-\allowbreak{}023, 027, 029 (exception: Llama Q4 448 control truncations in 024 and Llama v2 trailing period in 026). Disposition: Retained as canonical scoring standard through final release. \\[2pt]
\texttt{M09} & Gaussian-\allowbreak{}Copula Family-\allowbreak{}Cluster Power Simulation & 016 & 100,000 Monte Carlo replicates of paired binary differences with rho=0.50, discordance=0.70-\allowbreak{}0.90, two-\allowbreak{}sided family t-\allowbreak{}test. (Avoid underpowered pilots or treating correlated prompts within a family as independent N.) & Powered Cohort 016 (N=96 families, 80.7\% power for 20 pp effect) before collection. Disposition: Supplemented in Cohort 029 by stratified Welch-\allowbreak{}Satterthwaite coverage calibration (M19). \\[2pt]
\texttt{M10} & Signal Detection Theory (d', Criterion) \& Graded Strengths & 018, 019, 024, 025 & d' = Phi\textasciicircum{}-\allowbreak{}1(H) -\allowbreak{} Phi\textasciicircum{}-\allowbreak{}1(F), c = -\allowbreak{}0.5*[Phi\textasciicircum{}-\allowbreak{}1(H) + Phi\textasciicircum{}-\allowbreak{}1(F)] with half-\allowbreak{}count boundary correction on valid responses. (Balanced accuracy of 50\% conceals whether a model guesses randomly or activates 100\% of the time.) & Exposed negative/\allowbreak{}zero d' paired with extreme negative criterion (activation bias) in Qwen causal tasks vs positive d' in noncausal controls. Disposition: Retained as descriptive coordinates; two-\allowbreak{}point strength grid in 018 was insufficient for full psychometric curves. \\[2pt]
\texttt{M11} & Within-\allowbreak{}World Token-\allowbreak{}Mapping \& Option-\allowbreak{}Order Counterbalancing & 018, 019, 023, 024, 025, 027, 028, 029 & Semantic decoding (token XOR mapping) and metamorphic pair consistency checks. (Isolate superficial token/\allowbreak{}position habits from substantive evidence use.) & Uncovered +25.8/\allowbreak{}+36.7 pp K/\allowbreak{}M reversal shifts in Qwen (019) and 256/\allowbreak{}256 first-\allowbreak{}option selection in Mistral F16 (024/\allowbreak{}025). Disposition: Retained; crossed within-\allowbreak{}family in 023/\allowbreak{}024/\allowbreak{}027/\allowbreak{}028 and balanced across stratum slots in 029. \\[2pt]
\texttt{M12} & Staged Component \& Prompt-\allowbreak{}Guidance Decomposition & 020, 021, 022, 023, 024 & Family-\allowbreak{}paired differences (e.g., Acc(rand\_only) -\allowbreak{} Acc(two\_source)) in stateless fresh contexts. (Test whether two-\allowbreak{}source failure comes from misreading numbers or failing to recognize randomized designs.) & Demonstrated 100\% governing-\allowbreak{}source selection (128/\allowbreak{}128) alongside 45.3\% ordinary and guided action accuracy (58/\allowbreak{}128) in Qwen 7B F16. Disposition: Retained as behavioral dissociation evidence, with explicit caveat that separate calls do not prove internal causal mediation. \\[2pt]
\texttt{M13} & Matched Noncausal Scorecard \& Harmful-\allowbreak{}Direction Controls & 018, 019, 023, 024, 027, 028, 029 & Direction-\allowbreak{}stratified control accuracy (>=75\% in both beneficial and harmful directions) and causal-\allowbreak{}minus-\allowbreak{}control contrasts. (Prevent attributing general arithmetic, table-\allowbreak{}reading, or decline-\allowbreak{}avoidance errors to causal reasoning deficits.) & Blocked premature causal-\allowbreak{}specificity claims in 018, 023, 024, 027, and anchored the primary contrast in 029. Disposition: Retained as essential discriminant control through final manuscript. \\[2pt]
\texttt{M14} & Retrospective Metamorphic \& Information-\allowbreak{}Theoretic Archive Audit & 025, 025\_1 & I(Response; Truth) and I(Response; Assoc | Truth) over \{0, 1, invalid\}; paired item transition matrices. (Synthesize 24 cohorts (>40,000 responses) without spending compute or conflating post hoc checks with confirmation.) & Quantified Mistral's 6/\allowbreak{}128 option invariance, Qwen's 0.031-\allowbreak{}bit causal MI, and non-\allowbreak{}transfer of 016's Q4 order interaction to 024 (0 pp). Disposition: Completed its audit role and formally bounded non-\allowbreak{}estimable methods. \\[2pt]
\texttt{M15} & Exact-\allowbreak{}Rational Observationally Identical Graph-\allowbreak{}Switch Assay & 026, 027, 028 & Exact Python fractions.Fraction verification: tau\_conf = sum\_w P(w)[P(Y=1|X=1,w)-\allowbreak{}P(Y=1|X=0,w)] < 0 < tau\_med. (Eliminate observational table cues completely so only supplied causal structure determines the action.) & Mathematically airtight (after fixing v1 Y-\allowbreak{}complement bug in 026 and option-\allowbreak{}order coupling in 027), but all tested 3B-\allowbreak{}8B F16 models scored 0/\allowbreak{}8 order-\allowbreak{}robust paired switch fidelity. Disposition: Abandoned after Cohort 028 proved even calculator-\allowbreak{}assisted F16 models failed graph-\allowbreak{}to-\allowbreak{}rule selection. \\[2pt]
\texttt{M16} & Structured Rule Selection + Allowlisted Deterministic Calculator Replay & 028 & Primary gain mean\_i(C\_i -\allowbreak{} B\_i) on causal raw tasks vs specificity contrast (C-\allowbreak{}B)\_causal -\allowbreak{} (C-\allowbreak{}B)\_control. (Test whether arithmetic errors alone masked graph-\allowbreak{}identification competence in Cohort 027.) & Cleanly falsified the arithmetic-\allowbreak{}only bottleneck hypothesis: C rescued controls to 87.5\%-\allowbreak{}100\% while leaving causal accuracy near 50\%. Disposition: Completed in Cohort 028; motivated shift to identifiable study-\allowbreak{}provenance transport in Cohort 029. \\[2pt]
\texttt{M17} & Durable Append-\allowbreak{}Only Journal, Lockfiles \& Hourly/\allowbreak{}Branch Git Checkpoints & 001-\allowbreak{}029 (matured in 018, 024, 027, 028, 029) & Intent -\allowbreak{}> dispatch -\allowbreak{}> base64-\allowbreak{}encoded raw response envelope bound to manifest SHA-\allowbreak{}256. (Prevent silent replacement of failed/\allowbreak{}wrong attempts, survive multi-\allowbreak{}hour CPU pauses/\allowbreak{}shutdowns, and guarantee byte-\allowbreak{}identical offline replay.) & Enabled zero-\allowbreak{}loss recovery across Cohort 024 budget pauses, Cohort 028 shutdown, and 4,096/\allowbreak{}4,096 verified envelopes in Cohort 029. Disposition: Retained in final release. \\[2pt]
\texttt{M18} & CF-\allowbreak{}TP1 Feasible Role-\allowbreak{}Swap \& Target-\allowbreak{}Population Standardization & 029 & Latent surface (2h-\allowbreak{}r, 3r-\allowbreak{}2h) in (0,1)\textasciicircum{}8 for both roles; tau\_k = w\_k*delta\_1 + (1-\allowbreak{}w\_k)*delta\_0; pair fidelity F\_pair = I\_0 * I\_1. (Break Cohort 016's magnitude-\allowbreak{}gap shortcut and distinguish direction-\allowbreak{}correct switches from constant bias or random flipping.) & Zero public/\allowbreak{}private oracle mismatches across 4,096 prompts; exposed 995/\allowbreak{}1024 causal activation bias (3/\allowbreak{}256 switches) vs 45/\allowbreak{}256 unwarranted arithmetic flips. Disposition: Canonical methodology of the final research manuscript. \\[2pt]
\texttt{M19} & Dual-\allowbreak{}Oracle Verification, Enumerated Rival-\allowbreak{}Policy Gate \& Stratified Welch/\allowbreak{}Bounded Inference & 029 & SE = sqrt(sum\_h s\_h\textasciicircum{}2 /\allowbreak{} (64*16)), Welch df, plus distribution-\allowbreak{}free radius sqrt(2*log(40)/\allowbreak{}128) = +/\allowbreak{}-\allowbreak{}24.01 pp. (Guarantee prompt-\allowbreak{}to-\allowbreak{}key integrity, prove immunity to enumerated source-\allowbreak{}blind/\allowbreak{}target-\allowbreak{}blind shortcuts, and ensure valid interval coverage under discrete/\allowbreak{}bounded family scores.) & Passed 50/\allowbreak{}50 release verification checks; documented 12 literal constant-\allowbreak{}rival failures alongside 0 failures under constant-\allowbreak{}rivals-\allowbreak{}v1 clarification. Disposition: Canonical statistical and audit framework of the final manuscript. \\[2pt]
\end{longtable}
}

% =============================================================================
\chapter{Complete Research Decision Ledger}\label{app:decision_ledger}
Table~\ref{tab:decision_ledger} archives all thirty pivotal research and governance decisions (\texttt{D01}--\texttt{D30}) across the 320-commit Git history from \texttt{evidence/decision\_ledger.csv}.

{\footnotesize
\setlength{\tabcolsep}{3.5pt}
\begin{longtable}{@{}>{\raggedright\arraybackslash}p{0.9cm}>{\raggedright\arraybackslash}p{1.7cm}>{\raggedright\arraybackslash}p{1.3cm}>{\raggedright\arraybackslash}p{4.2cm}>{\raggedright\arraybackslash}p{4.0cm}>{\raggedright\arraybackslash}p{2.3cm}@{}}
\caption{Chronological Research Decision Ledger (\texttt{D01}--\texttt{D30}) with Git commits, documented rationales, and primary archive records.}\label{tab:decision_ledger}\\
\toprule
\textbf{ID} & \textbf{Date} & \textbf{Commit} & \textbf{Decision} & \textbf{Documented Rationale} & \textbf{Source Record} \\
\midrule
\endfirsthead
\toprule
\textbf{ID} & \textbf{Date} & \textbf{Commit} & \textbf{Decision} & \textbf{Documented Rationale} & \textbf{Source Record} \\
\midrule
\endhead
\bottomrule
\endfoot
\texttt{D01} & 2026-\allowbreak{}10-\allowbreak{}28 & \texttt{4e2fe4f} & Scope primary claim to behavioral causal competence rather than neural circuit severing & Accuracy/\allowbreak{}graph outputs cannot identify internal weight mechanisms without mechanistic interventions & \texttt{literature review/\allowbreak{}EXECUTIVE\_SUMMARY.md} \\[2pt]
\texttt{D02} & 2026-\allowbreak{}11-\allowbreak{}01 & \texttt{6ab18d2} & Separate public\_tasks.jsonl from evaluator-\allowbreak{}only oracle\_truth.jsonl with opaque SHA-\allowbreak{}256 task IDs & Prevent condition/\allowbreak{}answer leakage into LLM contexts & \texttt{docs/\allowbreak{}DATA\_CONTRACT.md} \\[2pt]
\texttt{D03} & 2026-\allowbreak{}11-\allowbreak{}07 & \texttt{5a12a21} & Enforce code-\allowbreak{}in-\allowbreak{}Git and weights-\allowbreak{}outside-\allowbreak{}Git storage policy; regenerate 7B Q4\_K\_M with b12176 & Ensure cross-\allowbreak{}machine reproducibility and uniform quantizer provenance & \texttt{docs/\allowbreak{}SSD\_ASSET\_REUSE.md} \\[2pt]
\texttt{D04} & 2026-\allowbreak{}11-\allowbreak{}10 13:10 & \texttt{aec9fd2} & Standardize on CPU-\allowbreak{}only llama.cpp b12176 (8 threads, F16 KV, temp 0, no prompt cache) & Lab RTX 4000 Ada used nouveau driver without working CUDA; common CPU runtime avoids backend confounds & \texttt{DECISIONS.md} \\[2pt]
\texttt{D05} & 2026-\allowbreak{}11-\allowbreak{}10 13:11 & \texttt{da4507b} & Retain strict JSON primary score and add separately labeled single-\allowbreak{}fence JSON sensitivity & 100\% of initial F16 outputs wrapped valid JSON in Markdown fences & \texttt{DECISIONS.md} \\[2pt]
\texttt{D06} & 2026-\allowbreak{}11-\allowbreak{}10 13:22 & \texttt{0333c67} & Add exploratory 3B Q4\_K\_M arm and run CPU repacking diagnostic on 3B Q2\_K & 3B Q2\_K hit 160-\allowbreak{}token cap with repetitive punctuation on 64/\allowbreak{}64 prompts & \texttt{DECISIONS.md} \\[2pt]
\texttt{D07} & 2026-\allowbreak{}11-\allowbreak{}10 13:38 & \texttt{941cf7e} & Retire 3B Q2\_K from subsequent pilots after -\allowbreak{}-\allowbreak{}no-\allowbreak{}repack and tokenizer audits & Both default and -\allowbreak{}-\allowbreak{}no-\allowbreak{}repack yielded 0/\allowbreak{}16 valid outputs despite identical tokenizers & \texttt{DECISIONS.md} \\[2pt]
\texttt{D08} & 2026-\allowbreak{}11-\allowbreak{}10 13:43 & \texttt{ea78a2f} & Normalize local home paths (/\allowbreak{}home/\allowbreak{}i242514/\allowbreak{}) to \$HOME/\allowbreak{} in tracked run artifacts without rewriting Git history & Protect lab account privacy while preserving immutable commit history & \texttt{DECISIONS.md} \\[2pt]
\texttt{D09} & 2026-\allowbreak{}11-\allowbreak{}10 13:50 & \texttt{f2b2cb1} & Score numeric ate\_estimate sign separately from explicit binary activate field & Cohort 002 F16 got 61/\allowbreak{}64 signs right while explicitly declining on 64/\allowbreak{}64 prompts & \texttt{DECISIONS.md} \\[2pt]
\texttt{D10} & 2026-\allowbreak{}11-\allowbreak{}10 14:07 & \texttt{01c7590} & Enforce Cohort 004 stopping gate and withhold 7B Q2\_K arm & 7B F16 scored 17/\allowbreak{}32 overall, missing predeclared 22/\allowbreak{}32 competence threshold & \texttt{DECISIONS.md} \\[2pt]
\texttt{D11} & 2026-\allowbreak{}11-\allowbreak{}10 14:13 & \texttt{4e90b24} & Withhold C/\allowbreak{}A/\allowbreak{}B latent-\allowbreak{}route parameter estimates and launch within-\allowbreak{}item source-\allowbreak{}order reversal & Cohort 005 F16 violated C/\allowbreak{}A/\allowbreak{}B additivity by -\allowbreak{}43.75 pp and showed family-\allowbreak{}confounded order split & \texttt{DECISIONS.md} \\[2pt]
\texttt{D12} & 2026-\allowbreak{}11-\allowbreak{}10 14:33 & \texttt{dd75c35} & Preserve interrupted Cohort 008 run\_001\_f16 (25/\allowbreak{}64 + 1 interrupted) and restart complete arms under new run IDs & User requested PC shutdown; partial unattempted items must not enter paired contrasts or be overwritten & \texttt{DECISIONS.md} \\[2pt]
\texttt{D13} & 2026-\allowbreak{}11-\allowbreak{}12 09:23 & \texttt{d21db20} & Introduce same-\allowbreak{}session interleaved order crossover with shuffled opaque twins (Cohorts 009-\allowbreak{}010) & Cohort 008 order contrast was confounded with a two-\allowbreak{}day session gap & \texttt{DECISIONS.md} \\[2pt]
\texttt{D14} & 2026-\allowbreak{}11-\allowbreak{}12 09:52 & \texttt{cd3f66d} & Replace fixed-\allowbreak{}confounding generator with matched observational-\allowbreak{}strength generator (Cohort 011) & Original generator made |delta\_obs| larger in aligned than conflict cells & \texttt{DECISIONS.md} \\[2pt]
\texttt{D15} & 2026-\allowbreak{}11-\allowbreak{}12 10:10 & \texttt{b90cb63} & Correct halved cell-\allowbreak{}count bug in analyze\_interleaved\_order.py and archive *\_v1\_undercounted\_counts.json & Code audit found family-\allowbreak{}mean counts were not multiplied by 2; point estimates/\allowbreak{}CIs unchanged & \texttt{DECISIONS.md} \\[2pt]
\texttt{D16} & 2026-\allowbreak{}11-\allowbreak{}12 10:21 & \texttt{fc7dcc1} & Require conflict competence separately in EACH evidence order (>=9/\allowbreak{}16) and withhold 7B Q4 in Cohort 012 & 7B F16 passed overall (47/\allowbreak{}64) via 32/\allowbreak{}32 aligned but failed both conflict orders (8/\allowbreak{}16, 7/\allowbreak{}16) & \texttt{DECISIONS.md} \\[2pt]
\texttt{D17} & 2026-\allowbreak{}11-\allowbreak{}12 11:36 & \texttt{a4a6f67} & Enforce strict bare one-\allowbreak{}key JSON instruction and 48-\allowbreak{}token cap (Cohort 014) & Cohort 013 causal-\allowbreak{}context prompts appended explanatory prose after JSON and hit 160-\allowbreak{}token cap on 64/\allowbreak{}64 items & \texttt{DECISIONS.md} \\[2pt]
\texttt{D18} & 2026-\allowbreak{}11-\allowbreak{}12 12:39 & \texttt{eb53578} & Freeze 96-\allowbreak{}family prospective replication of 7B Q4\_K\_M conflict-\allowbreak{}order effect (Cohort 016) & Cohort 015 revealed unexpected +65.6 pp Q4 conflict order shift; required untouched powered confirmation & \texttt{DECISIONS.md} \\[2pt]
\texttt{D19} & 2026-\allowbreak{}11-\allowbreak{}12 16:03 & \texttt{058e3c7} & Foreground causal inference in LLMs as overarching pillar; treat quantization as one experimental factor & Prevent narrowing project to prompt-\allowbreak{}order quirks or treating order sensitivity as causal reasoning & \texttt{DECISIONS.md} \\[2pt]
\texttt{D20} & 2026-\allowbreak{}11-\allowbreak{}13 & \texttt{caf5b1d} & Recover Cohort 017 14B paired analysis locally with disclosed zero-\allowbreak{}variance t-\allowbreak{}test guard; preserve original missing-\allowbreak{}NumPy log & Unattended lab script crashed post-\allowbreak{}inference due to missing NumPy and zero family variance in 14B Q4 & \texttt{DECISIONS.md} \\[2pt]
\texttt{D21} & 2026-\allowbreak{}11-\allowbreak{}14 & \texttt{a400569} & Issue dated interpretation amendment on Cohort 019 source-\allowbreak{}check position cue without altering scores or failed gate & Cohort 019 source check mentioned randomized study first in prose while asserting historical was listed first & \texttt{DECISIONS.md} \\[2pt]
\texttt{D22} & 2026-\allowbreak{}11-\allowbreak{}14 & \texttt{77361f5} & Stop single-\allowbreak{}checkpoint Qwen 7B prompt tuning after Cohort 021 guidance failed to improve 58/\allowbreak{}128 two-\allowbreak{}source decisions & 128/\allowbreak{}128 governing-\allowbreak{}source selection coexisted with 58/\allowbreak{}128 guided/\allowbreak{}ordinary decisions; needed cross-\allowbreak{}lineage check & \texttt{DECISIONS.md} \\[2pt]
\texttt{D23} & 2026-\allowbreak{}11-\allowbreak{}16 & \texttt{21dabb2} & Classify Qwen two-\allowbreak{}source dissociation as checkpoint-\allowbreak{}specific after Mistral failed simple components in Cohort 022 & Mistral Q4\_K\_M output constant 'A' (64/\allowbreak{}128) and constant 'activate' (16/\allowbreak{}32 rand-\allowbreak{}only) & \texttt{DECISIONS.md} \\[2pt]
\texttt{D24} & 2026-\allowbreak{}11-\allowbreak{}16 & \texttt{4aff668} & Withhold causal-\allowbreak{}specificity confirmation trigger after Cohort 023 harmful-\allowbreak{}direction controls failed (<75\%) & Qwen harmful numeric (18/\allowbreak{}32) and noncausal scorecard (75/\allowbreak{}128) missed frozen 75\% robustness cutoffs & \texttt{DECISIONS.md} \\[2pt]
\texttt{D25} & 2026-\allowbreak{}11-\allowbreak{}18 & \texttt{2fecac4} & Close 21-\allowbreak{}arm Cohort 024 as a development screen without confirmatory precision claims & No F16 reference cleared all component/\allowbreak{}control checks; precision patterns were heterogeneous and non-\allowbreak{}monotone & \texttt{DECISIONS.md} \\[2pt]
\texttt{D26} & 2026-\allowbreak{}11-\allowbreak{}19 & \texttt{3606794} & Pivot to exact-\allowbreak{}rational observationally identical confounder/\allowbreak{}mediator graph-\allowbreak{}switch assay (Cohort 026) & Cohorts 025/\allowbreak{}025.1 showed prior archives lacked fixed-\allowbreak{}table structural causal contrasts & \texttt{experiments/\allowbreak{}cohort\_025\_1/\allowbreak{}summary.md} \\[2pt]
\texttt{D27} & 2026-\allowbreak{}11-\allowbreak{}19 & \texttt{e118351} & Invalidate and preserve Cohort 026 Qwen v1 pilot due to Y-\allowbreak{}complement arithmetic-\allowbreak{}oracle bug; stop Cohort 026 after v3 gate failure & v1 hardcoded 'yes' on complemented outcomes; neither Qwen nor Llama passed v2/\allowbreak{}v3 qualification & \texttt{experiments/\allowbreak{}cohort\_026/\allowbreak{}development\_summary.md} \\[2pt]
\texttt{D28} & 2026-\allowbreak{}11-\allowbreak{}19 & \texttt{11a46f0 /\allowbreak{} 31fa454} & Exclude 14B F16 for 32 GB RAM guard, cross answer order structurally in Cohort 027, and honor stopped\_by\_gate & Cohort 026 coupled correct option position to form; all four F16 models scored 0/\allowbreak{}8 paired switch fidelity in 027 & \texttt{experiments/\allowbreak{}cohort\_027/\allowbreak{}README.md} \\[2pt]
\texttt{D29} & 2026-\allowbreak{}11-\allowbreak{}20 to 11-\allowbreak{}23 & \texttt{e7787e4 /\allowbreak{} 5a46c2e /\allowbreak{} ac6ed58} & Test deterministic calculator replay (Cohort 028), amend Git checkpoint cadence to hourly, and honor stopped\_by\_gate & Calculator rescued noncausal controls but not causal graph-\allowbreak{}rule selection (0/\allowbreak{}16 raw paired switch on 3 of 4 model-\allowbreak{}forms) & \texttt{experiments/\allowbreak{}cohort\_028/\allowbreak{}DEVIATIONS.md} \\[2pt]
\texttt{D30} & 2026-\allowbreak{}11-\allowbreak{}23 to 11-\allowbreak{}26 & \texttt{bf58d38 /\allowbreak{} 76206c2 /\allowbreak{} b737153} & Freeze and execute CF-\allowbreak{}TP1 (Cohort 029) with constant-\allowbreak{}rivals-\allowbreak{}v1 clarification, tag cohort-\allowbreak{}029-\allowbreak{}v1.0.0, and reconcile RESULTS.md direction table in manuscript & Provide shortcut-\allowbreak{}audited, offline-\allowbreak{}reproducible prospective evaluation of provenance and target transport with full provenance disclosure & \texttt{PUBLICATION\_HANDOFF.md} \\[2pt]
\end{longtable}
}
